\documentclass[journal,12pt,twocolumn]{IEEEtran}
%
\usepackage{setspace}
\usepackage{gensymb}
%\doublespacing
\singlespacing

%\usepackage{graphicx}
%\usepackage{amssymb}
%\usepackage{relsize}
\usepackage[cmex10]{amsmath}
%\usepackage{amsthm}
%\interdisplaylinepenalty=2500
%\savesymbol{iint}
%\usepackage{txfonts}
%\restoresymbol{TXF}{iint}
%\usepackage{wasysym}
\usepackage{amsthm}
\usepackage{iithtlc}
\usepackage{mathrsfs}
\usepackage{txfonts}
\usepackage{stfloats}
\usepackage{cite}
\usepackage{cases}
\usepackage{subfig}
%\usepackage{xtab}
\usepackage{longtable}
\usepackage{multirow}
%\usepackage{algorithm}
%\usepackage{algpseudocode}
\usepackage{enumitem}
\usepackage{mathtools}
\usepackage{tikz}
\usepackage{circuitikz}
\usepackage{verbatim}
\usepackage[breaklinks=true]{hyperref}
%\usepackage{stmaryrd}
\usepackage{tkz-euclide} % loads  TikZ and tkz-base
\usetkzobj{all}
\usepackage{listings}
    \usepackage{color}                                            %%
    \usepackage{array}                                            %%
    \usepackage{longtable}                                        %%
    \usepackage{calc}                                             %%
    \usepackage{multirow}                                         %%
    \usepackage{hhline}                                           %%
    \usepackage{ifthen}                                           %%
  %optionally (for landscape tables embedded in another document): %%
    \usepackage{lscape}     
\usepackage{multicol}
\usepackage{chngcntr}
%\usepackage{enumerate}

%\usepackage{wasysym}
%\newcounter{MYtempeqncnt}
\DeclareMathOperator*{\Res}{Res}
%\renewcommand{\baselinestretch}{2}
\renewcommand\thesection{\arabic{section}}
\renewcommand\thesubsection{\thesection.\arabic{subsection}}
\renewcommand\thesubsubsection{\thesubsection.\arabic{subsubsection}}

\renewcommand\thesectiondis{\arabic{section}}
\renewcommand\thesubsectiondis{\thesectiondis.\arabic{subsection}}
\renewcommand\thesubsubsectiondis{\thesubsectiondis.\arabic{subsubsection}}

% correct bad hyphenation here
\hyphenation{op-tical net-works semi-conduc-tor}
\def\inputGnumericTable{}                                 %%

\lstset{
%language=C,
frame=single, 
breaklines=true,
columns=fullflexible
}
%\lstset{
%language=tex,
%frame=single, 
%breaklines=true
%}

\begin{document}
%


\newtheorem{theorem}{Theorem}[section]
\newtheorem{problem}{Problem}
\newtheorem{proposition}{Proposition}[section]
\newtheorem{lemma}{Lemma}[section]
\newtheorem{corollary}[theorem]{Corollary}
\newtheorem{example}{Example}[section]
\newtheorem{definition}[problem]{Definition}
%\newtheorem{thm}{Theorem}[section] 
%\newtheorem{defn}[thm]{Definition}
%\newtheorem{algorithm}{Algorithm}[section]
%\newtheorem{cor}{Corollary}
\newcommand{\BEQA}{\begin{eqnarray}}
\newcommand{\EEQA}{\end{eqnarray}}
\newcommand{\define}{\stackrel{\triangle}{=}}

\bibliographystyle{IEEEtran}
%\bibliographystyle{ieeetr}

\providecommand{\nCr}[2]{\,^{#1}C_{#2}} % nCr
\providecommand{\nPr}[2]{\,^{#1}P_{#2}} % nPr

\providecommand{\mbf}{\mathbf}
\providecommand{\pr}[1]{\ensuremath{\Pr\left(#1\right)}}
\providecommand{\qfunc}[1]{\ensuremath{Q\left(#1\right)}}
\providecommand{\sbrak}[1]{\ensuremath{{}\left[#1\right]}}
\providecommand{\lsbrak}[1]{\ensuremath{{}\left[#1\right.}}
\providecommand{\rsbrak}[1]{\ensuremath{{}\left.#1\right]}}
\providecommand{\brak}[1]{\ensuremath{\left(#1\right)}}
\providecommand{\lbrak}[1]{\ensuremath{\left(#1\right.}}
\providecommand{\rbrak}[1]{\ensuremath{\left.#1\right)}}
\providecommand{\cbrak}[1]{\ensuremath{\left\{#1\right\}}}
\providecommand{\lcbrak}[1]{\ensuremath{\left\{#1\right.}}
\providecommand{\rcbrak}[1]{\ensuremath{\left.#1\right\}}}
\theoremstyle{remark}
\newtheorem{rem}{Remark}
\newcommand{\sgn}{\mathop{\mathrm{sgn}}}
\providecommand{\abs}[1]{\left\vert#1\right\vert}
\providecommand{\res}[1]{\Res\displaylimits_{#1}} 
\providecommand{\norm}[1]{\left\lVert#1\right\rVert}
%\providecommand{\norm}[1]{\lVert#1\rVert}
\providecommand{\mtx}[1]{\mathbf{#1}}
\providecommand{\mean}[1]{E\left[ #1 \right]}
\providecommand{\fourier}{\overset{\mathcal{F}}{ \rightleftharpoons}}
%\providecommand{\hilbert}{\overset{\mathcal{H}}{ \rightleftharpoons}}
\providecommand{\system}{\overset{\mathcal{H}}{ \longleftrightarrow}}
	%\newcommand{\solution}[2]{\textbf{Solution:}{#1}}
\newcommand{\solution}{\noindent \textbf{Solution: }}
\newcommand{\cosec}{\,\text{cosec}\,}
\providecommand{\dec}[2]{\ensuremath{\overset{#1}{\underset{#2}{\gtrless}}}}
\newcommand{\myvec}[1]{\ensuremath{\begin{pmatrix}#1\end{pmatrix}}}
\newcommand{\mydet}[1]{\ensuremath{\begin{vmatrix}#1\end{vmatrix}}}
%\numberwithin{equation}{section}
\numberwithin{equation}{subsection}
%\numberwithin{problem}{section}
%\numberwithin{definition}{section}
\makeatletter
\@addtoreset{figure}{problem}
\makeatother

\let\StandardTheFigure\thefigure
\let\vec\mathbf
%\renewcommand{\thefigure}{\theproblem.\arabic{figure}}
\renewcommand{\thefigure}{\theproblem}
%\setlist[enumerate,1]{before=\renewcommand\theequation{\theenumi.\arabic{equation}}
%\counterwithin{equation}{enumi}


%\renewcommand{\theequation}{\arabic{subsection}.\arabic{equation}}

\def\putbox#1#2#3{\makebox[0in][l]{\makebox[#1][l]{}\raisebox{\baselineskip}[0in][0in]{\raisebox{#2}[0in][0in]{#3}}}}
     \def\rightbox#1{\makebox[0in][r]{#1}}
     \def\centbox#1{\makebox[0in]{#1}}
     \def\topbox#1{\raisebox{-\baselineskip}[0in][0in]{#1}}
     \def\midbox#1{\raisebox{-0.5\baselineskip}[0in][0in]{#1}}

\vspace{3cm}

\title{
	\logo{
Discrete Mathematics
	}
}
\author{ G V V Sharma$^{*}$% <-this % stops a space
	\thanks{*The author is with the Department
		of Electrical Engineering, Indian Institute of Technology, Hyderabad
		502285 India e-mail:  gadepall@iith.ac.in. All content in this manual is released under GNU GPL.  Free and open source.}
	
}	
%\title{
%	\logo{Matrix Analysis through Octave}{\begin{center}\includegraphics[scale=.24]{tlc}\end{center}}{}{HAMDSP}
%}


% paper title
% can use linebreaks \\ within to get better formatting as desired
%\title{Matrix Analysis through Octave}
%
%
% author names and IEEE memberships
% note positions of commas and nonbreaking spaces ( ~ ) LaTeX will not break
% a structure at a ~ so this keeps an author's name from being broken across
% two lines.
% use \thanks{} to gain access to the first footnote area
% a separate \thanks must be used for each paragraph as LaTeX2e's \thanks
% was not built to handle multiple paragraphs
%

%\author{<-this % stops a space
%\thanks{}}
%}
% note the % following the last \IEEEmembership and also \thanks - 
% these prevent an unwanted space from occurring between the last author name
% and the end of the author line. i.e., if you had this:
% 
% \author{....lastname \thanks{...} \thanks{...} }
%                     ^------------^------------^----Do not want these spaces!
%
% a space would be appended to the last name and could cause every name on that
% line to be shifted left slightly. This is one of those "LaTeX things". For
% instance, "\textbf{A} \textbf{B}" will typeset as "A B" not "AB". To get
% "AB" then you have to do: "\textbf{A}\textbf{B}"
% \thanks is no different in this regard, so shield the last } of each \thanks
% that ends a line with a % and do not let a space in before the next \thanks.
% Spaces after \IEEEmembership other than the last one are OK (and needed) as
% you are supposed to have spaces between the names. For what it is worth,
% this is a minor point as most people would not even notice if the said evil
% space somehow managed to creep in.



% The paper headers
%\markboth{Journal of \LaTeX\ Class Files,~Vol.~6, No.~1, January~2007}%
%{Shell \MakeLowercase{\textit{et al.}}: Bare Demo of IEEEtran.cls for Journals}
% The only time the second header will appear is for the odd numbered pages
% after the title page when using the twoside option.
% 
% *** Note that you probably will NOT want to include the author's ***
% *** name in the headers of peer review papers.                   ***
% You can use \ifCLASSOPTIONpeerreview for conditional compilation here if
% you desire.




% If you want to put a publisher's ID mark on the page you can do it like
% this:
%\IEEEpubid{0000--0000/00\$00.00~\copyright~2007 IEEE}
% Remember, if you use this you must call \IEEEpubidadjcol in the second
% column for its text to clear the IEEEpubid mark.



% make the title area
\maketitle

\newpage

\tableofcontents

\bigskip

\renewcommand{\thefigure}{\theenumi}
\renewcommand{\thetable}{\theenumi}
%\renewcommand{\theequation}{\theenumi}

%\begin{abstract}
%%\boldmath
%In this letter, an algorithm for evaluating the exact analytical bit error rate  (BER)  for the piecewise linear (PL) combiner for  multiple relays is presented. Previous results were available only for upto three relays. The algorithm is unique in the sense that  the actual mathematical expressions, that are prohibitively large, need not be explicitly obtained. The diversity gain due to multiple relays is shown through plots of the analytical BER, well supported by simulations. 
%
%\end{abstract}
% IEEEtran.cls defaults to using nonbold math in the Abstract.
% This preserves the distinction between vectors and scalars. However,
% if the journal you are submitting to favors bold math in the abstract,
% then you can use LaTeX's standard command \boldmath at the very start
% of the abstract to achieve this. Many IEEE journals frown on math
% in the abstract anyway.

% Note that keywords are not normally used for peerreview papers.
%\begin{IEEEkeywords}
%Cooperative diversity, decode and forward, piecewise linear
%\end{IEEEkeywords}



% For peer review papers, you can put extra information on the cover
% page as needed:
% \ifCLASSOPTIONpeerreview
% \begin{center} \bfseries EDICS Category: 3-BBND \end{center}
% \fi
%
% For peerreview papers, this IEEEtran command inserts a page break and
% creates the second title. It will be ignored for other modes.
%\IEEEpeerreviewmaketitle

\begin{abstract}
This book provides a computational approach to discrete mathematics by solving problems in related areas  from IIT-JEE.  Links to sample C/Python codes are available in the text.  The book provides sufficient math basics for Machine Learning and is also recommended for high school students who wish to explore topics in Artificial Intelligence.  
\end{abstract}
Download python codes using 
\begin{lstlisting}
svn co https://github.com/gadepall/school/trunk/jee/discrete/codes
\end{lstlisting}
%\section{JEE Exercises}

\documentclass[journal,12pt,twocolumn]{IEEEtran}
\usepackage{setspace}
\usepackage{gensymb}
\usepackage{caption}
%\usepackage{multirow}
%\usepackage{multicolumn}
%\usepackage{subcaption}
%\doublespacing
\singlespacing
\usepackage{csvsimple}
\usepackage{amsmath}
\usepackage{multicol}
%\usepackage{enumerate}
\usepackage{amssymb}
%\usepackage{graphicx}
\usepackage{newfloat}
%\usepackage{syntax}
\usepackage{listings}
\usepackage{iithtlc}
\usepackage{color}
\usepackage{tikz}
\usetikzlibrary{shapes,arrows}



%\usepackage{graphicx}
%\usepackage{amssymb}
%\usepackage{relsize}
%\usepackage[cmex10]{amsmath}
%\usepackage{mathtools}
%\usepackage{amsthm}
%\interdisplaylinepenalty=2500
%\savesymbol{iint}
%\usepackage{txfonts}
%\restoresymbol{TXF}{iint}
%\usepackage{wasysym}
\usepackage{amsthm}
\usepackage{mathrsfs}
\usepackage{txfonts}
\usepackage{stfloats}
\usepackage{cite}
\usepackage{cases}
\usepackage{mathtools}
\usepackage{caption}
\usepackage{enumerate}	
\usepackage{enumitem}
\usepackage{amsmath}
%\usepackage{xtab}
\usepackage{longtable}
\usepackage{multirow}
%\usepackage{algorithm}
%\usepackage{algpseudocode}
\usepackage{enumitem}
\usepackage{mathtools}
\usepackage{hyperref}
%\usepackage[framemethod=tikz]{mdframed}
\usepackage{listings}
    %\usepackage[latin1]{inputenc}                                 %%
    \usepackage{color}                                            %%
    \usepackage{array}                                            %%
    \usepackage{longtable}                                        %%
    \usepackage{calc}                                             %%
    \usepackage{multirow}                                         %%
    \usepackage{hhline}                                           %%
    \usepackage{ifthen}                                           %%
  %optionally (for landscape tables embedded in another document): %%
    \usepackage{lscape}     


\usepackage{url}
\def\UrlBreaks{\do\/\do-}


%\usepackage{stmaryrd}


%\usepackage{wasysym}
%\newcounter{MYtempeqncnt}
\DeclareMathOperator*{\Res}{Res}
%\renewcommand{\baselinestretch}{2}
\renewcommand\thesection{\arabic{section}}
\renewcommand\thesubsection{\thesection.\arabic{subsection}}
\renewcommand\thesubsubsection{\thesubsection.\arabic{subsubsection}}

\renewcommand\thesectiondis{\arabic{section}}
\renewcommand\thesubsectiondis{\thesectiondis.\arabic{subsection}}
\renewcommand\thesubsubsectiondis{\thesubsectiondis.\arabic{subsubsection}}

% correct bad hyphenation here
\hyphenation{op-tical net-works semi-conduc-tor}

%\lstset{
%language=C,
%frame=single, 
%breaklines=true
%}

%\lstset{
	%%basicstyle=\small\ttfamily\bfseries,
	%%numberstyle=\small\ttfamily,
	%language=Octave,
	%backgroundcolor=\color{white},
	%%frame=single,
	%%keywordstyle=\bfseries,
	%%breaklines=true,
	%%showstringspaces=false,
	%%xleftmargin=-10mm,
	%%aboveskip=-1mm,
	%%belowskip=0mm
%}

%\surroundwithmdframed[width=\columnwidth]{lstlisting}
\def\inputGnumericTable{}                                 %%
\lstset{
%language=C,
frame=single, 
breaklines=true,
columns=fullflexible
}
 

\begin{document}
%
\tikzstyle{block} = [rectangle, draw,
    text width=3em, text centered, minimum height=3em]
\tikzstyle{sum} = [draw, circle, node distance=3cm]
\tikzstyle{input} = [coordinate]
\tikzstyle{output} = [coordinate]
\tikzstyle{pinstyle} = [pin edge={to-,thin,black}]

\theoremstyle{definition}
\newtheorem{theorem}{Theorem}[section]
\newtheorem{problem}{Problem}
\newtheorem{proposition}{Proposition}[section]
\newtheorem{lemma}{Lemma}[section]
\newtheorem{corollary}[theorem]{Corollary}
\newtheorem{example}{Example}[section]
\newtheorem{definition}{Definition}[section]
%\newtheorem{algorithm}{Algorithm}[section]
%\newtheorem{cor}{Corollary}
\newcommand{\BEQA}{\begin{eqnarray}}
\newcommand{\EEQA}{\end{eqnarray}}
\newcommand{\define}{\stackrel{\triangle}{=}}

\bibliographystyle{IEEEtran}
%\bibliographystyle{ieeetr}

\providecommand{\nCr}[2]{\,^{#1}C_{#2}} % nCr
\providecommand{\nPr}[2]{\,^{#1}P_{#2}} % nPr
\providecommand{\mbf}{\mathbf}
\providecommand{\pr}[1]{\ensuremath{\Pr\left(#1\right)}}
\providecommand{\qfunc}[1]{\ensuremath{Q\left(#1\right)}}
\providecommand{\sbrak}[1]{\ensuremath{{}\left[#1\right]}}
\providecommand{\lsbrak}[1]{\ensuremath{{}\left[#1\right.}}
\providecommand{\rsbrak}[1]{\ensuremath{{}\left.#1\right]}}
\providecommand{\brak}[1]{\ensuremath{\left(#1\right)}}
\providecommand{\lbrak}[1]{\ensuremath{\left(#1\right.}}
\providecommand{\rbrak}[1]{\ensuremath{\left.#1\right)}}
\providecommand{\cbrak}[1]{\ensuremath{\left\{#1\right\}}}
\providecommand{\lcbrak}[1]{\ensuremath{\left\{#1\right.}}
\providecommand{\rcbrak}[1]{\ensuremath{\left.#1\right\}}}
\theoremstyle{remark}
\newtheorem{rem}{Remark}
\newcommand{\sgn}{\mathop{\mathrm{sgn}}}
\providecommand{\abs}[1]{\left\vert#1\right\vert}
\providecommand{\res}[1]{\Res\displaylimits_{#1}} 
\providecommand{\norm}[1]{\left\Vert#1\right\Vert}
\providecommand{\mtx}[1]{\mathbf{#1}}
\providecommand{\mean}[1]{E\left[ #1 \right]}
\providecommand{\fourier}{\overset{\mathcal{F}}{ \rightleftharpoons}}
%\providecommand{\hilbert}{\overset{\mathcal{H}}{ \rightleftharpoons}}
\providecommand{\system}{\overset{\mathcal{H}}{ \longleftrightarrow}}
	%\newcommand{\solution}[2]{\textbf{Solution:}{#1}}
\newcommand{\solution}{\noindent \textbf{Solution: }}
\newcommand{\myvec}[1]{\ensuremath{\begin{pmatrix}#1\end{pmatrix}}}
\providecommand{\dec}[2]{\ensuremath{\overset{#1}{\underset{#2}{\gtrless}}}}
\DeclarePairedDelimiter{\ceil}{\lceil}{\rceil}
%\numberwithin{equation}{subsection}
%\numberwithin{equation}{section}
%\numberwithin{problem}{subsection}
%\numberwithin{definition}{subsection}
\makeatletter
\@addtoreset{figure}{section}
\makeatother

\let\StandardTheFigure\thefigure
%\renewcommand{\thefigure}{\theproblem.\arabic{figure}}
\renewcommand{\thefigure}{\thesection}


%\numberwithin{figure}{subsection}

%\numberwithin{equation}{subsection}
%\numberwithin{equation}{section}
%\numberwithin{equation}{problem}
%\numberwithin{problem}{subsection}
\numberwithin{problem}{section}
%%\numberwithin{definition}{subsection}
%\makeatletter
%\@addtoreset{figure}{problem}
%\makeatother
\makeatletter
\@addtoreset{table}{section}
\makeatother

\let\StandardTheFigure\thefigure
\let\StandardTheTable\thetable
\let\vec\mathbf
%%\renewcommand{\thefigure}{\theproblem.\arabic{figure}}
%\renewcommand{\thefigure}{\theproblem}

%%\numberwithin{figure}{section}

%%\numberwithin{figure}{subsection}



\def\putbox#1#2#3{\makebox[0in][l]{\makebox[#1][l]{}\raisebox{\baselineskip}[0in][0in]{\raisebox{#2}[0in][0in]{#3}}}}
     \def\rightbox#1{\makebox[0in][r]{#1}}
     \def\centbox#1{\makebox[0in]{#1}}
     \def\topbox#1{\raisebox{-\baselineskip}[0in][0in]{#1}}
     \def\midbox#1{\raisebox{-0.5\baselineskip}[0in][0in]{#1}}

\vspace{3cm}

\title{ 
	\logo{
JEE Problems in Discrete Mathematics
	}
}

%\author{ G V V Sharma$^{*}$% <-this % stops a space
%	\thanks{*The author is with the Department
%		of Electrical Engineering, Indian Institute of Technology, Hyderabad
%		502285 India e-mail:  gadepall@iith.ac.in. All content in this manual is released under GNU GPL.  Free and open source.}
	
%}	

\maketitle

%\tableofcontents

\bigskip

\renewcommand{\thefigure}{\theenumi}
\renewcommand{\thetable}{\theenumi}


\begin{abstract}
	A  collection of problems from JEE papers related to matrices are available in this document.  Verify your soluions using  C.
\end{abstract}
\section{Signal Processing: Z Transform}
\begin{enumerate}[label=\thesection.\arabic*
,ref=\thesection.\theenumi]

\item Let
\begin{align}
a(n) &= \frac{\alpha^n-\beta^n}{\alpha-\beta}u(n)
\\
b(n) &=a(n-1)+a(n+1) - \delta(n)
\end{align}
where $\alpha, \beta$ are the roots of the equation
\begin{equation}
z^2-z-1 = 0
\end{equation}
and
\begin{align}
u(n)
=
\begin{cases}
0, & n < 0
\\
1, & n \ge 0
\end{cases}
\\
\delta(n)
=
\begin{cases}
0, & n \ne 0
\\
1, & n = 0
\end{cases}
\end{align}
\item Verify your results through a C program.
%\begin{enumerate}[label=\theenumi.\arabic*
%,ref=\thesection.\theenumi]
\item Show that the $Z$ transform of $u(n)$
\begin{align}
U(z)&\triangleq \sum_{n=-\infty}^{\infty}u(n)z^{-n} 
\\
&= \frac{1}{1- z^{-1}}, \quad \abs{z} > 1
\end{align}
\item  Show that 
\begin{align}
A(z)&= \frac{z^{-1}}{1- z^{-1}-z^{-2}}
\end{align}
\item Let 
\begin{align}
y(n)=a(n)*u(n) \triangleq \sum_{k=-\infty}^{\infty}a(k) u(n-k)
\end{align}
%
Show that 
\begin{align}
y(n) = \sum_{k=0}^{n}a(k) 
\end{align}
%
\item Show that 
\begin{align}
Y(z) &= A(z)U(z)
\\
&= \frac{z^{-1}}{\brak{1- z^{-1}-z^{-2}}\brak{1- z^{-1}}}
\end{align}
\item Show that 
\begin{align}
w(n)&=\sbrak{a(n+2)-1}u(n-1)
\\
&=a(n+2)-u(n+1)+2\delta(n)
\end{align}
\item Is $W(z)=Y(z)$?
\item Verify if 
\begin{align}
\sum_{n=1}^{\infty}\frac{a(n)}{10^n} = \frac{10}{89}
\end{align}
\item Verify if 
\begin{align}
\sum_{n=1}^{\infty}\frac{b(n)}{10^n} = \frac{8}{89}
\end{align}
\section{Algebra: Modular Arithmetic}
Let $AP\brak{a;d}$ denote an A.P. with $d > 0$
\begin{enumerate}[label=\thesection.\arabic*
,ref=\thesection.\theenumi]
\item Express $AP\brak{a;d}$ in modulo arithmetic.
\\
\solution 
\begin{align}
\label{eq:14_ap_mod_def}
A &\equiv a\pmod{d}
\end{align}
\item Express the intersection of $AP\brak{1;3}, AP\brak{2;5}$ and $AP\brak{3;7}$  using modulo arithmetic.
\\
\solution The desired AP can be expressed as
\begin{align}
\label{eq:14_ap_mod}
A &\equiv 1\pmod{3}
\\
& \equiv 2\pmod {5}
\\
&\equiv 3\pmod{7}
\end{align}
%
\item Two numbers are said to be coprime if their greatest common divisor (gcd) is 1. Verify if (3,5), (5,7) and (3,7) are pairwise coprime.
\item Does a solution for \eqref{eq:14_ap_mod} exist?
\\
\solution  The Chinese remainder theorem guarantees that the system in \eqref{eq:14_ap_mod} has a solution since 3,5,7 are pairwise coprime.
%\item Find the modular multiplicative inverse of 
%\begin{align}
%\label{eq:14_ap_modinv}
%A &\equiv \brak{7\times 11}\pmod{5}
%\end{align}
% 
\item Simplify 
\begin{align}
\label{eq:14_ap_mod_753}
\brak{7\times 5}\pmod{3}
\end{align}
%
\solution \eqref{eq:14_ap_mod_753} can be expressed as
%
\begin{align}
\label{eq:14_ap_mod_753_sol}
\brak{7\times 5}\pmod{3} &= 35\pmod{3} 
\nonumber \\
&=2\pmod{3}
\end{align}

\item Find $x$ in 
%
\begin{align}
\label{eq:14_ap_mod_753_prob}
2x =1\pmod{3}
\end{align}
%
%
\solution By inspection, for $x = 2$, 
\begin{align}
%\label{eq:14_ap_mod_753_sol}
2x = 2\times 2 = 4 = 3+1 =1\pmod{3}
\end{align}
%
Thus $x = 2$ is a solution of \eqref{eq:14_ap_mod_753_prob}. 

\item In general, $x$ in 
\begin{align}
\label{eq:14_ap_mod_inv_def}
ax =1\pmod{d}
\end{align}
%
is defined to be the modular multiplicative inverse of \eqref{eq:14_ap_mod_def}.
\item Show that the multiplicative inverse of 
\begin{align}
\label{eq:14_ap_mod_357_sol}
\brak{3\times 5}\pmod{7} &= y=1
\end{align}
%
\item Show that the multiplicative inverse of 
\begin{align}
\label{eq:14_ap_mod_375_sol}
\brak{3\times 7}\pmod{5} &= z=1
\end{align}
%
\item Find $a+d$.
\\
\solution 
\begin{multline}
\label{eq:14_ap_mod_ad_sol}
\brak{5\times 7 \times 1 \times x}+
\brak{3\times 5\times 3 \times y} 
\\
+
\brak{3\times 7\times 2 \times z} =  157
\end{multline}
%
\item Find $a$ and $d$.
\\
\solution
\begin{align}
\label{eq:14_ap_mod_d}
d &= LCM\brak{3,5,7} = 105
\\
A&=157 \pmod{105} 
\nonumber \\
&= 52\pmod{105} 
\nonumber \\
\implies a &= 52
\end{align}
\item Given the APs 
\begin{align}
\label{eq:14_ap_mod_gen}
a_1&\pmod{d_1}
\\
a_2 &\pmod {d_2}
\\
a_3 &\pmod {d_3},
\end{align}
%
such that 
\begin{align}
gcd(d_1,d_2)=gcd(d_2,d_3)=gcd(d_3,d_1)=1,
\end{align}
show that their intersection
\begin{align}
\label{eq:14_ap_mod_gen_int}
a\pmod{d}
\end{align}
%
is obtained through
\begin{multline}
\label{eq:14_ap_mod_ad_gen}
a+d =
\\
 \brak{d_1\times d_2 \times a_3 \times x}+
\brak{d_2\times d_3 \times a_1 \times y}
\\
+
\brak{d_3\times d_1 \times a_2 \times z}
\end{multline}
\begin{align}
\label{eq:14_ap_mod_d_gen}
d &= LCM\brak{d_1,d_2,d_3},
\end{align}
%
where $x, y, z$ are the modular multiplicative inverses given by
\begin{align}
\label{eq:14_ap_mod_mul_inv_gen}
x &= \sbrak{ \brak{d_1\times d_2}\pmod{d_3}}^{-1}
\\
y &= \sbrak{ \brak{d_2\times d_3}\pmod{d_1}}^{-1}
\\
z &= \sbrak{  \brak{d_3\times d_1}\pmod{d_2}}^{-1}
\end{align}
respectively.
%
\item Write  a C program to find $x,y$ and $z$.
%\item Find the vector $\myvec{n \\ k \\ r}$ with the smallest length such that $A_n \cap B_k \cap C_r $.
%\\
%\solution This is obtained by considering the following set of equations
%\begin{align}
% 3n-2= 5k-3
%\\
%5k-3 = 7r-4
%%\\
%%7r-4 = a + md
%\end{align}
%%
%which can be simplified to obtain the matrix equation 
%\begin{align}
%\myvec{3 & -5 & 0 \\ 0 & 5 & -7  }\myvec{n \\ k \\ r} = \myvec{-1 \\ -1 }
%%\myvec{3 & 5 & 0 \\ 0 & 5 & 7 \\ 7 &  0 & 0 }\myvec{n \\ k \\ r} = \myvec{-1 \\ -1 \\ a + md}
%\end{align}
%%
%yielding the augmented matrix 
%\begin{align}
%\myvec{3 & -5 & 0 & -1\\ 0 & 5 & -7 &-1 }\leftrightarrow \myvec{3 & 0 & -7 & -2\\ 0 & 5 & -7 &-1 } 
%\\
%\implies \myvec{3n \\ 5k \\ r} = \myvec{ -2 \\ -1\\ 0 } + r \myvec{ 7 \\ 7  \\ 1} 
%\\
%\implies \myvec{n \\ k \\ r} = \myvec{ \frac{7r-2}{3} \\ \frac{7r-1}{5}\\ r }
%%\myvec{3 & 5 & 0 \\ 0 & 5 & 7 \\ 7 &  0 & 0 }\myvec{n \\ k \\ r} = \myvec{-1 \\ -1 \\ a + md}
%\end{align}
%%
%For $r = 8, n = 18, k = 11$ are integers.  
%\item  Find $a+d$ such that $A_n \cap B_k \cap C_r \cap AP\brak{a:d}$.
%\\
%\solution Substituting $r = 8$ in \eqref{eq:14_r},
%\begin{align}
%a = 7r - 4 = 52
%\end{align}
%Also, 
%\begin{align}
%a+d &= a + 3i = a + 5j = a + 7m,
%\\
%\implies d &= 3i=5j=7m
%\end{align}
%%
%for some positive integers $i, j, m$. The smallest integer $d$ satisfying the above condition is $LCM\brak{3,5,7} = 105$.  Hence, $a + d = 52+105 = 157$.
\end{enumerate}

\end{enumerate}
\section{Discrete Fourier Transform}
\begin{enumerate}[label=\thesection.\arabic*
,ref=\thesection.\theenumi]
\item Show that 
\begin{align}
\label{eq:qp2_3_cexp}
\sum_{k=0}^{n-1}e^{\j \frac{2 \pi k}{n}} 
= 
\begin{cases}
1 & n = 1,
\\
0 & n > 1
\end{cases}
\end{align}
\item Show that 
%
\begin{align}
\label{eq:qp2_3_cosr}
\sum_{k=0}^{n}\cos \brak{\frac{2k+r}{n+2}\pi} = -\cos\brak{\frac{r-2}{n+2}\pi} 
\end{align}
\solution From \eqref{eq:qp2_3_cexp},
\begin{align}
\label{eq:qp2_3_cosr_sol}
\sum_{k=0}^{n+1}e^{\j \frac{2k+r}{n+2}\pi} &=0
\nonumber \\
\implies \sum_{k=0}^{n}e^{\j \frac{2k+r}{n+2}\pi} + e^{\j \frac{2\brak{n+1}+r}{n+2}\pi} &=0
\nonumber \\
\implies \sum_{k=0}^{n}e^{\j \frac{2k+r}{n+2}\pi} &=  -e^{\j \frac{2\brak{n+2}+r-2}{n+2}\pi} 
\nonumber \\
&=-e^{\j \frac{r-2}{n+2}\pi} 
\end{align}
Taking the real part on both sides yields \eqref{eq:qp2_3_cosr}.
%
\item Show that 
\begin{align}
\label{eq:qp2_3_sum}
f(n) & = \frac{\sum_{k=0}^{n}\sin\brak{\frac{k+1}{n+2}\pi}
\sin\brak{\frac{k+2}{n+2}\pi} }{\sum_{k=0}^{n}\sin^2\brak{\frac{k+2}{n+2}\pi}}
%\nonumber 
\\
& = 
\frac{\brak{n+1}\cos \brak{\frac{\pi}{n+2}}}{n+\cos\brak{\frac{2\pi}{n+2}}}
\label{eq:qp2_3_sum_ans}
\end{align}
%\frac{n\cos \cbrak{\frac{\pi}{n+2}}-\sum_{k=0}^{n}\cos \cbrak{\frac{2k+3}{n+2}\pi}}{n-\sum_{k=0}^{n}\cos\cbrak{\frac{2k+4}{n+2}\pi}}
\solution 
Let
\begin{align}
\label{eq:qp2_3_theta}
\theta_n = \frac{\pi}{n+2}
\end{align}
%
\begin{multline}
\label{eq:qp2_3_cos}
\because \sin\cbrak{\brak{k+1}\theta_n}
\sin\cbrak{\brak{k+2}\theta_n} ,
\\
= \frac{1}{2}\sbrak{\cos \theta_n-\cos \cbrak{\brak{2k+3}\theta_n}}
\end{multline}
from \eqref{eq:qp2_3_sum} and \eqref{eq:qp2_3_cosr},
\begin{align}
\label{eq:qp2_3_sum_sol}
f(n) & = \frac{n\cos \theta_n-\sum_{k=0}^{n}\cos \cbrak{\brak{2k+3}\theta_n}}{n-\sum_{k=0}^{n}\cos\cbrak{\brak{2k+4}\theta_n}}
\nonumber \\
&= \frac{n\cos \brak{\frac{\pi}{n+2}}+\cos \brak{\frac{\pi}{n+2}}}{n+\cos\brak{\frac{2\pi}{n+2}}}
\end{align}
resulting in \eqref{eq:qp2_3_sum_ans}.  Verify if 
\item 
\begin{align}
f(4) = \frac{\sqrt{3}}{2}
\end{align}
\item 
\begin{align}
\lim_{n \to \infty}f(n) = \frac{1}{2}
\end{align}
\item 
\begin{align}
\sin\brak{7 \cos^{-1}f(5)} = 0
\end{align}
\item If 
\begin{align}
\alpha = 
\tan\brak{ \cos^{-1}f(6)} 
\end{align}
verify if 
\begin{align}
\alpha^2 +2\alpha -1 =  0 
\end{align}
\end{enumerate}
%
\section{Combinatorics}
\begin{enumerate}[label=\thesection.\arabic*
,ref=\thesection.\theenumi]
\item Find 
\begin{align}
\label{eq:2019_qp2_comb_nat}
\sum_{k=0}^{n}k
\end{align}
\solution  \eqref{eq:2019_qp2_comb_nat} can be expressed as
\begin{align}
\label{eq:2019_qp2_comb_nat_sol}
\frac{n(n+1)}{2}
\end{align}
\item Find 
\begin{align}
\label{eq:2019_qp2_comb_k2}
\sum_{k=0}^{n}\nCr{n}{k}k^2
\end{align}
\solution 
\begin{align}
\label{eq:2019_qp2_comb}
\brak{1+x}^n&=\sum_{k=0}^{n}\nCr{n}{k}x^k
%\nonumber 
\\
\implies n\brak{1+x}^{n-1}&=\sum_{k=0}^{n}k\nCr{n}{k}x^{k-1}
\label{eq:2019_qp2_comb_k2_bin}
\end{align}
upon differentiation. Multiplying \eqref{eq:2019_qp2_comb_k2_bin} by $x$ and differentiating,
\begin{align}
\frac{d}{dx} \sbrak{nx\brak{1+x}^{n-1}}&=\sum_{k=0}^{n}k^2\nCr{n}{k}x^{k-1}
\end{align}
\begin{multline}
\implies n\brak{n-1}x\brak{1+x}^{n-2}+n\brak{1+x}^{n-1}  
\\
=\sum_{k=0}^{n}k^2\nCr{n}{k}x^{k-1}
\label{eq:2019_qp2_comb_k2_bin_diff}
\end{multline}
%
Substituting $x=1$ in \eqref{eq:2019_qp2_comb_k2_bin_diff},
\begin{align}
\sum_{k=0}^{n}\nCr{n}{k}k^2 &= n\brak{n-1}2^{n-2}+n2^{n-1}
\nonumber \\
&= n\brak{n+1}2^{n-2}
\label{eq:2019_qp2_comb_k2_bin_sol}
\end{align}
\item Find 
\begin{align}
\label{eq:2019_qp2_comb_k}
\sum_{k=0}^{n}\nCr{n}{k}k
\end{align}
\solution Substituting $x=1$ in \eqref{eq:2019_qp2_comb_k2_bin},
\begin{align}
\label{eq:2019_qp2_comb_k_sol} 
\sum_{k=0}^{n}\nCr{n}{k}k= n2^{n-1}
\end{align}
\item Find 
\begin{align}
\label{eq:2019_qp2_comb_3k}
\sum_{k=0}^{n}\nCr{n}{k}3^k
\end{align}
\solution Substituting $x=2$ in \eqref{eq:2019_qp2_comb},
\begin{align}
\label{eq:2019_qp2_comb_3k_sol}
\sum_{k=0}^{n}\nCr{n}{k}3^k = 4^n
\end{align}
\item If
\begin{align}
\label{eq:2019_qp2_comb_det}
\begin{vmatrix}
\frac{n(n+1)}{2} & n\brak{n+1}2^{n-2}
\\
n2^{n-1} & 4^n
\end{vmatrix} = 0
\end{align}
for some $n$, find 
\begin{align}
\label{eq:2019_qp2_comb_det_prob}
\sum_{k=0}^{n}\frac{\nCr{n}{k}}{k+1}
\end{align}
%
\solution \eqref{eq:2019_qp2_comb_det} can be expressed as
\begin{align}
\label{eq:2019_qp2_comb_det_sol}
n(n+1)2^{2n-3}
\begin{vmatrix}
1 & 1
\\
n & 4
\end{vmatrix} = 0
\\
\implies n = 4
\end{align}
%
Integrating \eqref{eq:2019_qp2_comb} from 0 to 1,

\begin{align}
\label{eq:2019_qp2_comb_int}
\frac{2^{n+1}}{n+1}&=\sum_{k=0}^{n}\frac{\nCr{n}{k}}{k+1}
\end{align}
%
Substituting $n=4$ in the above, 
\begin{align}
\label{eq:2019_qp2_comb_int_sol}
\sum_{k=0}^{n}\frac{\nCr{n}{k}}{k+1}= \frac{2^{5}-1}{5} = \frac{31}{5}
\end{align}
\end{enumerate}
\section{Probability}
Table \ref{table:2019_7} lists the number of red (R) and green (G) balls in bags $B_1, B_2$ and $B_3$. Also listed are the probabilities of each bag.
\begin{table}[!h]
\centering
%\resizebox {0.5\columnwidth} {!} {
%%%%%%%%%%%%%%%%%%%%%%%%%%%%%%%%%%%%%%%%%%%%%%%%%%%%%%%%%%%%%%%%%%%%%%%
%%                                                                  %%
%%  This is the header of a LaTeX2e file exported from Gnumeric.    %%
%%                                                                  %%
%%  This file can be compiled as it stands or included in another   %%
%%  LaTeX document. The table is based on the longtable package so  %%
%%  the longtable options (headers, footers...) can be set in the   %%
%%  preamble section below (see PRAMBLE).                           %%
%%                                                                  %%
%%  To include the file in another, the following two lines must be %%
%%  in the including file:                                          %%
%%        \def\inputGnumericTable{}                                 %%
%%  at the beginning of the file and:                               %%
%%        \input{name-of-this-file.tex}                             %%
%%  where the table is to be placed. Note also that the including   %%
%%  file must use the following packages for the table to be        %%
%%  rendered correctly:                                             %%
%%    \usepackage[latin1]{inputenc}                                 %%
%%    \usepackage{color}                                            %%
%%    \usepackage{array}                                            %%
%%    \usepackage{longtable}                                        %%
%%    \usepackage{calc}                                             %%
%%    \usepackage{multirow}                                         %%
%%    \usepackage{hhline}                                           %%
%%    \usepackage{ifthen}                                           %%
%%  optionally (for landscape tables embedded in another document): %%
%%    \usepackage{lscape}                                           %%
%%                                                                  %%
%%%%%%%%%%%%%%%%%%%%%%%%%%%%%%%%%%%%%%%%%%%%%%%%%%%%%%%%%%%%%%%%%%%%%%



%%  This section checks if we are begin input into another file or  %%
%%  the file will be compiled alone. First use a macro taken from   %%
%%  the TeXbook ex 7.7 (suggestion of Han-Wen Nienhuys).            %%
\def\ifundefined#1{\expandafter\ifx\csname#1\endcsname\relax}


%%  Check for the \def token for inputed files. If it is not        %%
%%  defined, the file will be processed as a standalone and the     %%
%%  preamble will be used.                                          %%
\ifundefined{inputGnumericTable}

%%  We must be able to close or not the document at the end.        %%
	\def\gnumericTableEnd{\end{document}}


%%%%%%%%%%%%%%%%%%%%%%%%%%%%%%%%%%%%%%%%%%%%%%%%%%%%%%%%%%%%%%%%%%%%%%
%%                                                                  %%
%%  This is the PREAMBLE. Change these values to get the right      %%
%%  paper size and other niceties.                                  %%
%%                                                                  %%
%%%%%%%%%%%%%%%%%%%%%%%%%%%%%%%%%%%%%%%%%%%%%%%%%%%%%%%%%%%%%%%%%%%%%%

	\documentclass[12pt%
			  %,landscape%
                    ]{report}
       \usepackage[latin1]{inputenc}
       \usepackage{fullpage}
       \usepackage{color}
       \usepackage{array}
       \usepackage{longtable}
       \usepackage{calc}
       \usepackage{multirow}
       \usepackage{hhline}
       \usepackage{ifthen}

	\begin{document}


%%  End of the preamble for the standalone. The next section is for %%
%%  documents which are included into other LaTeX2e files.          %%
\else

%%  We are not a stand alone document. For a regular table, we will %%
%%  have no preamble and only define the closing to mean nothing.   %%
    \def\gnumericTableEnd{}

%%  If we want landscape mode in an embedded document, comment out  %%
%%  the line above and uncomment the two below. The table will      %%
%%  begin on a new page and run in landscape mode.                  %%
%       \def\gnumericTableEnd{\end{landscape}}
%       \begin{landscape}


%%  End of the else clause for this file being \input.              %%
\fi

%%%%%%%%%%%%%%%%%%%%%%%%%%%%%%%%%%%%%%%%%%%%%%%%%%%%%%%%%%%%%%%%%%%%%%
%%                                                                  %%
%%  The rest is the gnumeric table, except for the closing          %%
%%  statement. Changes below will alter the table's appearance.     %%
%%                                                                  %%
%%%%%%%%%%%%%%%%%%%%%%%%%%%%%%%%%%%%%%%%%%%%%%%%%%%%%%%%%%%%%%%%%%%%%%

\providecommand{\gnumericmathit}[1]{#1} 
%%  Uncomment the next line if you would like your numbers to be in %%
%%  italics if they are italizised in the gnumeric table.           %%
%\renewcommand{\gnumericmathit}[1]{\mathit{#1}}
\providecommand{\gnumericPB}[1]%
{\let\gnumericTemp=\\#1\let\\=\gnumericTemp\hspace{0pt}}
 \ifundefined{gnumericTableWidthDefined}
        \newlength{\gnumericTableWidth}
        \newlength{\gnumericTableWidthComplete}
        \newlength{\gnumericMultiRowLength}
        \global\def\gnumericTableWidthDefined{}
 \fi
%% The following setting protects this code from babel shorthands.  %%
 \ifthenelse{\isundefined{\languageshorthands}}{}{\languageshorthands{english}}
%%  The default table format retains the relative column widths of  %%
%%  gnumeric. They can easily be changed to c, r or l. In that case %%
%%  you may want to comment out the next line and uncomment the one %%
%%  thereafter                                                      %%
\providecommand\gnumbox{\makebox[0pt]}
%%\providecommand\gnumbox[1][]{\makebox}

%% to adjust positions in multirow situations                       %%
\setlength{\bigstrutjot}{\jot}
\setlength{\extrarowheight}{\doublerulesep}

%%  The \setlongtables command keeps column widths the same across  %%
%%  pages. Simply comment out next line for varying column widths.  %%
\setlongtables

\setlength\gnumericTableWidth{%
	133pt+%
	53pt+%
	57pt+%
0pt}
\def\gumericNumCols{3}
\setlength\gnumericTableWidthComplete{\gnumericTableWidth+%
         \tabcolsep*\gumericNumCols*2+\arrayrulewidth*\gumericNumCols}
\ifthenelse{\lengthtest{\gnumericTableWidthComplete > \linewidth}}%
         {\def\gnumericScale{\ratio{\linewidth-%
                        \tabcolsep*\gumericNumCols*2-%
                        \arrayrulewidth*\gumericNumCols}%
{\gnumericTableWidth}}}%
{\def\gnumericScale{1}}

%%%%%%%%%%%%%%%%%%%%%%%%%%%%%%%%%%%%%%%%%%%%%%%%%%%%%%%%%%%%%%%%%%%%%%
%%                                                                  %%
%% The following are the widths of the various columns. We are      %%
%% defining them here because then they are easier to change.       %%
%% Depending on the cell formats we may use them more than once.    %%
%%                                                                  %%
%%%%%%%%%%%%%%%%%%%%%%%%%%%%%%%%%%%%%%%%%%%%%%%%%%%%%%%%%%%%%%%%%%%%%%

\ifthenelse{\isundefined{\gnumericColA}}{\newlength{\gnumericColA}}{}\settowidth{\gnumericColA}{\begin{tabular}{@{}p{100pt*\gnumericScale}@{}}x\end{tabular}}
\ifthenelse{\isundefined{\gnumericColB}}{\newlength{\gnumericColB}}{}\settowidth{\gnumericColB}{\begin{tabular}{@{}p{50pt*\gnumericScale}@{}}x\end{tabular}}
\ifthenelse{\isundefined{\gnumericColC}}{\newlength{\gnumericColC}}{}\settowidth{\gnumericColC}{\begin{tabular}{@{}p{70pt*\gnumericScale}@{}}x\end{tabular}}

%\begin{longtable}[c]{%
\begin{table}[!h]


\begin{tabular}[c]{
	b{\gnumericColA}%
	b{\gnumericColB}%
	b{\gnumericColC}%
	}

%%%%%%%%%%%%%%%%%%%%%%%%%%%%%%%%%%%%%%%%%%%%%%%%%%%%%%%%%%%%%%%%%%%%%%
%%  The longtable options. (Caption, headers... see Goosens, p.124) %%
%	\caption{The Table Caption.}             \\	%
% \hline	% Across the top of the table.
%%  The rest of these options are table rows which are placed on    %%
%%  the first, last or every page. Use \multicolumn if you want.    %%

%%  Header for the first page.                                      %%
%	\multicolumn{3}{c}{The First Header} \\ \hline 
%	\multicolumn{1}{c}{colTag}	%Column 1
%	&\multicolumn{1}{c}{colTag}	%Column 2
%	&\multicolumn{1}{c}{colTag}	\\ \hline %Last column
%	\endfirsthead

%%  The running header definition.                                  %%
%	\hline
%	\multicolumn{3}{l}{\ldots\small\slshape continued} \\ \hline
%	\multicolumn{1}{c}{colTag}	%Column 1
%	&\multicolumn{1}{c}{colTag}	%Column 2
%	&\multicolumn{1}{c}{colTag}	\\ \hline %Last column
%	\endhead

%%  The running footer definition.                                  %%
%	\hline
%	\multicolumn{3}{r}{\small\slshape continued\ldots} \\
%	\endfoot

%%  The ending footer definition.                                   %%
%	\multicolumn{3}{c}{That's all folks} \\ \hline 
%	\endlastfoot
%%%%%%%%%%%%%%%%%%%%%%%%%%%%%%%%%%%%%%%%%%%%%%%%%%%%%%%%%%%%%%%%%%%%%%

\hhline{|-|-|-}
	 \multicolumn{1}{|p{\gnumericColA}|}%
	{\gnumericPB{\centering}\textbf{Component}}
	&\multicolumn{1}{p{\gnumericColB}|}%
	{\gnumericPB{\raggedright}\textbf{Value}}
	&\multicolumn{1}{p{\gnumericColC}|}%
	{\gnumericPB{\centering}\textbf{Quantity}}
\\
\hhline{|---|}
	 \multicolumn{1}{|p{\gnumericColA}|}%
	{\gnumericPB{\centering}Breadboard}
	&\multicolumn{1}{p{\gnumericColB}|}%
	{\gnumericPB{\raggedright} }
	&\multicolumn{1}{p{\gnumericColC}|}%
	{\gnumericPB{\centering}1}
\\
\hhline{|---|}
	 \multicolumn{1}{|p{\gnumericColA}|}%
	{\gnumericPB{\centering}Resistor}
	&\multicolumn{1}{p{\gnumericColB}|}%
	{\gnumericPB{\raggedright} $\ge 220 \Omega$}
	&\multicolumn{1}{p{\gnumericColC}|}%
	{\gnumericPB{\centering}1}
\\
\hhline{|---|}
	 \multicolumn{1}{|p{\gnumericColA}|}%
	{\gnumericPB{\centering}Pi}
	&\multicolumn{1}{p{\gnumericColB}|}%
	{Model B, Rev 3}
	&\multicolumn{1}{p{\gnumericColC}|}%
	{\gnumericPB{\centering}1}
\\
\hhline{|---|}
	 \multicolumn{1}{|p{\gnumericColA}|}%
	{\gnumericPB{\centering}Seven Segment Display}
	&\multicolumn{1}{p{\gnumericColB}|}%
	{Common Anode}
	&\multicolumn{1}{p{\gnumericColC}|}%
	{\gnumericPB{\centering}1}
\\
\hhline{|---|}
	 \multicolumn{1}{|p{\gnumericColA}|}%
	{\gnumericPB{\centering}\gnumbox{Jumper Wires}}
	&\multicolumn{1}{p{\gnumericColB}|}%
	{Female-Male}
	&\multicolumn{1}{p{\gnumericColC}|}%
	{\gnumericPB{\centering}\gnumbox{20}}
\\
\hhline{|-|-|-|}
%\end{longtable}
\end{tabular}
\caption{}
\label{table:components}
\end{table}
\ifthenelse{\isundefined{\languageshorthands}}{}{\languageshorthands{\languagename}}
\gnumericTableEnd

%%%%%%%%%%%%%%%%%%%%%%%%%%%%%%%%%%%%%%%%%%%%%%%%%%%%%%%%%%%%%%%%%%%%%%
%%                                                                  %%
%%  This is the header of a LaTeX2e file exported from Gnumeric.    %%
%%                                                                  %%
%%  This file can be compiled as it stands or included in another   %%
%%  LaTeX document. The table is based on the longtable package so  %%
%%  the longtable options (headers, footers...) can be set in the   %%
%%  preamble section below (see PRAMBLE).                           %%
%%                                                                  %%
%%  To include the file in another, the following two lines must be %%
%%  in the including file:                                          %%
%%        \def\inputGnumericTable{}                                 %%
%%  at the beginning of the file and:                               %%
%%        \input{name-of-this-file.tex}                             %%
%%  where the table is to be placed. Note also that the including   %%
%%  file must use the following packages for the table to be        %%
%%  rendered correctly:                                             %%
%%    \usepackage[latin1]{inputenc}                                 %%
%%    \usepackage{color}                                            %%
%%    \usepackage{array}                                            %%
%%    \usepackage{longtable}                                        %%
%%    \usepackage{calc}                                             %%
%%    \usepackage{multirow}                                         %%
%%    \usepackage{hhline}                                           %%
%%    \usepackage{ifthen}                                           %%
%%  optionally (for landscape tables embedded in another document): %%
%%    \usepackage{lscape}                                           %%
%%                                                                  %%
%%%%%%%%%%%%%%%%%%%%%%%%%%%%%%%%%%%%%%%%%%%%%%%%%%%%%%%%%%%%%%%%%%%%%%



%%  This section checks if we are begin input into another file or  %%
%%  the file will be compiled alone. First use a macro taken from   %%
%%  the TeXbook ex 7.7 (suggestion of Han-Wen Nienhuys).            %%
\def\ifundefined#1{\expandafter\ifx\csname#1\endcsname\relax}


%%  Check for the \def token for inputed files. If it is not        %%
%%  defined, the file will be processed as a standalone and the     %%
%%  preamble will be used.                                          %%
\ifundefined{inputGnumericTable}

%%  We must be able to close or not the document at the end.        %%
	\def\gnumericTableEnd{\end{document}}


%%%%%%%%%%%%%%%%%%%%%%%%%%%%%%%%%%%%%%%%%%%%%%%%%%%%%%%%%%%%%%%%%%%%%%
%%                                                                  %%
%%  This is the PREAMBLE. Change these values to get the right      %%
%%  paper size and other niceties.                                  %%
%%                                                                  %%
%%%%%%%%%%%%%%%%%%%%%%%%%%%%%%%%%%%%%%%%%%%%%%%%%%%%%%%%%%%%%%%%%%%%%%

	\documentclass[12pt%
			  %,landscape%
                    ]{report}
       \usepackage[latin1]{inputenc}
       \usepackage{fullpage}
       \usepackage{color}
       \usepackage{array}
       \usepackage{longtable}
       \usepackage{calc}
       \usepackage{multirow}
       \usepackage{hhline}
       \usepackage{ifthen}

	\begin{document}


%%  End of the preamble for the standalone. The next section is for %%
%%  documents which are included into other LaTeX2e files.          %%
\else

%%  We are not a stand alone document. For a regular table, we will %%
%%  have no preamble and only define the closing to mean nothing.   %%
    \def\gnumericTableEnd{}

%%  If we want landscape mode in an embedded document, comment out  %%
%%  the line above and uncomment the two below. The table will      %%
%%  begin on a new page and run in landscape mode.                  %%
%       \def\gnumericTableEnd{\end{landscape}}
%       \begin{landscape}


%%  End of the else clause for this file being \input.              %%
\fi

%%%%%%%%%%%%%%%%%%%%%%%%%%%%%%%%%%%%%%%%%%%%%%%%%%%%%%%%%%%%%%%%%%%%%%
%%                                                                  %%
%%  The rest is the gnumeric table, except for the closing          %%
%%  statement. Changes below will alter the table's appearance.     %%
%%                                                                  %%
%%%%%%%%%%%%%%%%%%%%%%%%%%%%%%%%%%%%%%%%%%%%%%%%%%%%%%%%%%%%%%%%%%%%%%

\providecommand{\gnumericmathit}[1]{#1} 
%%  Uncomment the next line if you would like your numbers to be in %%
%%  italics if they are italizised in the gnumeric table.           %%
%\renewcommand{\gnumericmathit}[1]{\mathit{#1}}
\providecommand{\gnumericPB}[1]%
{\let\gnumericTemp=\\#1\let\\=\gnumericTemp\hspace{0pt}}
 \ifundefined{gnumericTableWidthDefined}
        \newlength{\gnumericTableWidth}
        \newlength{\gnumericTableWidthComplete}
        \newlength{\gnumericMultiRowLength}
        \global\def\gnumericTableWidthDefined{}
 \fi
%% The following setting protects this code from babel shorthands.  %%
 \ifthenelse{\isundefined{\languageshorthands}}{}{\languageshorthands{english}}
%%  The default table format retains the relative column widths of  %%
%%  gnumeric. They can easily be changed to c, r or l. In that case %%
%%  you may want to comment out the next line and uncomment the one %%
%%  thereafter                                                      %%
\providecommand\gnumbox{\makebox[0pt]}
%%\providecommand\gnumbox[1][]{\makebox}

%% to adjust positions in multirow situations                       %%
\setlength{\bigstrutjot}{\jot}
\setlength{\extrarowheight}{\doublerulesep}

%%  The \setlongtables command keeps column widths the same across  %%
%%  pages. Simply comment out next line for varying column widths.  %%
\setlongtables

\setlength\gnumericTableWidth{%
	28pt+%
	12pt+%
	13pt+%
	75pt+%
0pt}
\def\gumericNumCols{4}
\setlength\gnumericTableWidthComplete{\gnumericTableWidth+%
         \tabcolsep*\gumericNumCols*2+\arrayrulewidth*\gumericNumCols}
\ifthenelse{\lengthtest{\gnumericTableWidthComplete > \linewidth}}%
         {\def\gnumericScale{\ratio{\linewidth-%
                        \tabcolsep*\gumericNumCols*2-%
                        \arrayrulewidth*\gumericNumCols}%
{\gnumericTableWidth}}}%
{\def\gnumericScale{1}}

%%%%%%%%%%%%%%%%%%%%%%%%%%%%%%%%%%%%%%%%%%%%%%%%%%%%%%%%%%%%%%%%%%%%%%
%%                                                                  %%
%% The following are the widths of the various columns. We are      %%
%% defining them here because then they are easier to change.       %%
%% Depending on the cell formats we may use them more than once.    %%
%%                                                                  %%
%%%%%%%%%%%%%%%%%%%%%%%%%%%%%%%%%%%%%%%%%%%%%%%%%%%%%%%%%%%%%%%%%%%%%%

\ifthenelse{\isundefined{\gnumericColA}}{\newlength{\gnumericColA}}{}\settowidth{\gnumericColA}{\begin{tabular}{@{}p{28pt*\gnumericScale}@{}}x\end{tabular}}
\ifthenelse{\isundefined{\gnumericColB}}{\newlength{\gnumericColB}}{}\settowidth{\gnumericColB}{\begin{tabular}{@{}p{12pt*\gnumericScale}@{}}x\end{tabular}}
\ifthenelse{\isundefined{\gnumericColC}}{\newlength{\gnumericColC}}{}\settowidth{\gnumericColC}{\begin{tabular}{@{}p{13pt*\gnumericScale}@{}}x\end{tabular}}
\ifthenelse{\isundefined{\gnumericColD}}{\newlength{\gnumericColD}}{}\settowidth{\gnumericColD}{\begin{tabular}{@{}p{75pt*\gnumericScale}@{}}x\end{tabular}}

\begin{tabular}[c]{%
	b{\gnumericColA}%
	b{\gnumericColB}%
	b{\gnumericColC}%
	b{\gnumericColD}%
	}

%%%%%%%%%%%%%%%%%%%%%%%%%%%%%%%%%%%%%%%%%%%%%%%%%%%%%%%%%%%%%%%%%%%%%%
%%  The longtable options. (Caption, headers... see Goosens, p.124) %%
%	\caption{The Table Caption.}             \\	%
% \hline	% Across the top of the table.
%%  The rest of these options are table rows which are placed on    %%
%%  the first, last or every page. Use \multicolumn if you want.    %%

%%  Header for the first page.                                      %%
%	\multicolumn{4}{c}{The First Header} \\ \hline 
%	\multicolumn{1}{c}{colTag}	%Column 1
%	&\multicolumn{1}{c}{colTag}	%Column 2
%	&\multicolumn{1}{c}{colTag}	%Column 3
%	&\multicolumn{1}{c}{colTag}	\\ \hline %Last column
%	\endfirsthead

%%  The running header definition.                                  %%
%	\hline
%	\multicolumn{4}{l}{\ldots\small\slshape continued} \\ \hline
%	\multicolumn{1}{c}{colTag}	%Column 1
%	&\multicolumn{1}{c}{colTag}	%Column 2
%	&\multicolumn{1}{c}{colTag}	%Column 3
%	&\multicolumn{1}{c}{colTag}	\\ \hline %Last column
%	\endhead

%%  The running footer definition.                                  %%
%	\hline
%	\multicolumn{4}{r}{\small\slshape continued\ldots} \\
%	\endfoot

%%  The ending footer definition.                                   %%
%	\multicolumn{4}{c}{That's all folks} \\ \hline 
%	\endlastfoot
%%%%%%%%%%%%%%%%%%%%%%%%%%%%%%%%%%%%%%%%%%%%%%%%%%%%%%%%%%%%%%%%%%%%%%

\hhline{|-|-|-|-}
	 \multicolumn{1}{|p{\gnumericColA}|}%
	{\gnumericPB{\centering}\gnumbox{\textbf{Bag}}}
	&\multicolumn{1}{p{\gnumericColB}|}%
	{\gnumericPB{\centering}\gnumbox{\textbf{R}}}
	&\multicolumn{1}{p{\gnumericColC}|}%
	{\gnumericPB{\centering}\gnumbox{\textbf{G}}}
	&\multicolumn{1}{p{\gnumericColD}|}%
	{\gnumericPB{\centering}\gnumbox{\textbf{Probability}}}
\\
\hhline{|----|}
	 \multicolumn{1}{|p{\gnumericColA}|}%
	{\gnumericPB{\centering}\gnumbox{$B_1$}}
	&\multicolumn{1}{p{\gnumericColB}|}%
	{\gnumericPB{\centering}\gnumbox{5}}
	&\multicolumn{1}{p{\gnumericColC}|}%
	{\gnumericPB{\centering}\gnumbox{5}}
	&\multicolumn{1}{p{\gnumericColD}|}%
	{\gnumericPB{\centering}\gnumbox{$\Pr(B_1)=\frac{3}{10}$}}
\\
\hhline{|----|}
	 \multicolumn{1}{|p{\gnumericColA}|}%
	{\gnumericPB{\centering}\gnumbox{$B_2$}}
	&\multicolumn{1}{p{\gnumericColB}|}%
	{\gnumericPB{\centering}\gnumbox{3}}
	&\multicolumn{1}{p{\gnumericColC}|}%
	{\gnumericPB{\centering}\gnumbox{5}}
	&\multicolumn{1}{p{\gnumericColD}|}%
	{\gnumericPB{\centering}\gnumbox{$\Pr(B_2)=\frac{3}{10}$}}
\\
\hhline{|----|}
	 \multicolumn{1}{|p{\gnumericColA}|}%
	{\gnumericPB{\centering}\gnumbox{$B_3$}}
	&\multicolumn{1}{p{\gnumericColB}|}%
	{\gnumericPB{\centering}\gnumbox{5}}
	&\multicolumn{1}{p{\gnumericColC}|}%
	{\gnumericPB{\centering}\gnumbox{3}}
	&\multicolumn{1}{p{\gnumericColD}|}%
	{\gnumericPB{\centering}\gnumbox{$\Pr(B_3)=\frac{4}{10}$}}
\\
\hhline{|-|-|-|-|}
\end{tabular}

\ifthenelse{\isundefined{\languageshorthands}}{}{\languageshorthands{\languagename}}
\gnumericTableEnd

%}
\caption{}
\label{table:2019_7}
\end{table}

\begin{enumerate}[label=\thesection.\arabic*
,ref=\thesection.\theenumi]
\item Show that 
\begin{align}
\pr{G|B_3} = \frac{3}{8}
\end{align}
\item Show that 
\begin{align}
\pr{G} = \frac{39}{80} 
\end{align}
\solution 
\begin{align}
\because \pr{G|B_1} &= \frac{1}{2}, \pr{G|B_2} = \frac{5}{8}, \pr{G|B_3} = \frac{3}{8},
\nonumber \\
\pr{G} &= \sum_{i=1}^{3}\pr{G|B_i}\pr{B_i}
\\
 &= \frac{1}{2}\times\frac{3}{10}+\frac{5}{8} \times\frac{3}{10}+\frac{3}{8}\times \frac{4}{10}
\\
 &= \frac{39}{80}
\end{align}
\item Is
\begin{align}
\pr{B_3|G} = \frac{5}{13}\, ?
\end{align}
\solution 
\begin{align}
\pr{B_3|G} &= \frac{\pr{G|B_3}\pr{B_3}}{\pr{G}}
\\
&= \frac{\frac{3}{8}\times \frac{4}{10}}{\frac{39}{80}} = \frac{4}{13} \ne \frac{5}{13}
\end{align}
\item Is
\begin{align}
\pr{B_3 \cap G} = \frac{3}{10}\, ?
\end{align}
\solution 
\begin{align}
\pr{B_3 \cap G} &= \pr{G|B_3}\pr{B_3} \\
\\
=& \frac{3}{8} \times \frac{4}{10}=\frac{3}{20}\ne\frac{3}{10}
\end{align}
\end{enumerate}
\section{Combinatorics}
\begin{enumerate}[label=\thesection.\arabic*
,ref=\thesection.\theenumi]
\item Five persons $A, B, C, D $ and $E$ are seated in a circuar arrangment. Let $A, C, E$ be given the colour green. Find the number of ways to distribute blue and red to $B$ and $D$.
\\
\solution Both $B$  and $D$ can be given either blue or red.  The number of possible ways is
\begin{align}
\label{eq:qp2_10_ace_g}
2\times 2 = 4
\end{align}
\item Repeat the above exercise with $B,D, A$ having the colour green.
\item Find the total number of ways in which the colour green can be distributed to alternate persons so that persons seated in adjacent seats get different coloured hats.
\\
\solution The number of such ways is
\begin{align}
\label{eq:qp2_10_g}
2\times 2 \times 2= 8
\end{align}
\item  If each person  is given a hat of one of 3 colours red, blue and green, then find the number of ways of distributing the hats such that the persons seated in adjacent seats get different coloured hats.
\\
\solution The number of ways is 
\begin{align}
\label{eq:qp2_10_sol}
2\times 2 \times 2 \times 3= 24
\end{align}
\end{enumerate}
\end{document}

\section{JEE Exercises:Sequences and Series}
\renewcommand{\theequation}{\theenumi}
\begin{enumerate}[label=\arabic*.,ref=\thesubsection.\theenumi]
\numberwithin{equation}{enumi}
    \item The sum of integers from 1 to 100 that are divisible by 2 or 5 is....................... 
     \item The solution of the equation $log_7 log_5(\sqrt{x+5}+\sqrt{x})=0$ is ....................
     \item The sum of the first n terms of the series $1^2+2.2^2+3^2+2.4^2+5^2+2.6^2+................$is $\frac{n(n+1)^2}{2}$, when n is even.When n is odd,the sum is ...............
     \item Let the harmonic mean and geometric mean of two positive numbers be the ratio 4:5.Then the two numbers are in the ratio...................
     \item For any odd integer n$\geq1$,$n^3-(n-1)^3+...+(-1)^{(n-1)} 1^3$=.............
     \item Let p and q be roots of the equation 
     \begin{align}
     x^2-2x+A=0
     \end{align} and let r and s be the roots of the equation 
     \begin{align}
     x^2-18x+B=0
     \end{align}. If $p<q<r<s$ are in arithmetic progression,then A =.......... and B=........

 \textbf{ MCQs with One Correct Answer}
 
 \item If x,y and z are $p^th$,$q^th$ and $r^th$ terms respectively of an A.P.and also of a G.P.,then $x^y-z y^z-x z^x-y$ is equal to:
 \begin{enumerate}
 \item xyz 
 \item 0 
 \item 1 
 \item None of these
 \end{enumerate}
\item The third term of a geometric progression is 4.The product of the first five terms is 
\begin{enumerate}
 \item $4^3$ 
 \item $4^5$
 \item $4^4$
 \item None of these
 \end{enumerate}
 \item The rational number,which equals the number $2.\overline{357}$ with recurring decimal is
 \begin{enumerate}
 \item $\frac{2355}{1001}$ 
\item  $\frac{2379}{997}$ 
 \item $\frac{2355}{999}$
 \item none of these
 \end{enumerate}
 \item If a,b,c are in G.P.,then the equations
  \begin{align}
  ax^2+2bx+c=0
  \end{align} and  
  \begin{align}
  dx^2+2ex+f=0
  \end{align} have a common root if $\frac{d}{a}, \frac{e}{b}, \frac{f}{c}$ are in......
  \begin{enumerate}
  \item A.P. 
  \item G.P. 
  \item H.P. 
  \item None of these
 \end{enumerate}
  \item Sum of the first n terms of the series $\frac{1}{2}$+$\frac{3}{4}$+$\frac{7}{8}$+$\frac{15}{16}$+.....is equal to 
   \begin{enumerate}
 \item  $2^n-n-1$ 
 \item $1-2^-n$
 \item  $n+2^-n-1$
 \item $2^n+1$
  \end{enumerate}
\item The number $log_2$7 is
   \begin{enumerate}
\item an integer 
\item a rational number 
\item an irrational number  
\item a prime number
   \end{enumerate}
  \item If $ln(a+c)$, $ln(a-c)$, $ln(a-2b+c)$ are in A.P.,then
 \begin{enumerate} 
\item a,b,c are in A.P. 
\item $a^2$, $b^2$, $c^2$ are in A.P.
\item a,b,c are in G.P. 
\item a,b,c are in H.P.
 \end{enumerate}
  \item Let $a_1,a_2,.....a_10$ be in A.P.and $h_1,h_2,.....h_{10}$ be in H.P.If $a_1=h_1=2$ and $a_{10}=h_{10}$=3, then $a_4 h_7$ is
 \begin{enumerate}
\item 2
\item 3 
\item 5 
\item 6
\end{enumerate}
  \item The harmonic mean of the roots of the equation 
  \begin{align}
  (5+\sqrt{2})x^2- (4+\sqrt{5})x+ 8 + 2\sqrt{5}=0
  \end{align} is
\begin{enumerate}
\item 2
\item 4 
\item 6 
\item 8
\end{enumerate}
\item Consider an infinite geometric series with first term a and common ratio r. If its sum is 4 and the second term is $\frac{3}{4}$, then
\begin{enumerate}
\item a=$\frac{4}{7}$,r=$\frac{3}{7}$ 
\item a=2,r=$\frac{3}{8}$
\item a=$\frac{3}{2}$,r=$\frac{1}{2}$
\item a=3,r=$\frac{1}{4}$
\end{enumerate}
\item Let $\alpha$, $\beta$ be the roots of 
\begin{align}
x^2-x+p=0
\end{align} and $\gamma$, $\delta$ be the roots of 
\begin{align}
x^2-4x+q=0
\end{align}. If $\alpha$, $\beta$, $\gamma$, $\delta$ are in G.P., then the integral values of p and q respectively, are
\begin{enumerate}
\item -2,-32 
\item -2,3 
\item -6,3 
\item -6,-32
\end{enumerate}
\item Let the positive numbers a,b,c,d be in A.P.Then abc,abd,acd,bcd are
\begin{enumerate}
\item NOT in A.P./G.P./H.P. 
\item in A.P.
\item in G.P.
\item in H.P.
\end{enumerate}
\item If the sum of the first 2n terms of the A.P.2,5,8,.....,is equal to the sum of the first n terms of the A.P.57,59,61,.....,then n equals
\begin{enumerate}
\item 10 
\item 12 
\item 11 
\item 13
\end{enumerate}
\item Suppose a,b,c are in A.P. and $a^2, b^2, c^2$ are in G.P.if $a<b<c$ and a+b+c=$\frac{3}{2}$, then the value of a is
\begin{enumerate}
\item $\frac{1}{2\sqrt{2}}$ 
\item $\frac{1}{2\sqrt{3}}$ 
\item $\frac{1}{2}$-$\frac{1}{\sqrt{3}}$ 
\item $\frac{1}{2}$-$\frac{1}{\sqrt{2}}$
\end{enumerate}
\item An infinite G.P. has first term 'x' and sum '5', then x belongs to
\begin{enumerate}
\item $x<-10$ 
\item $-10<x<0$ 
\item $0<x<10$ 
\item $x>10$
\end{enumerate}
\item In the quadratic equation 
\begin{align}
ax^2+bx+c=0,
\end{align} $\Delta =b^2-4ac$ and $\alpha+\beta$, $\alpha^2+\beta^2$, $\alpha^3+\beta^3$, are in G.P.where $\alpha, \beta$ are the root of $ax^2+bx+c=0$, then 
\begin{enumerate}
\item $\Delta\neq0$
\item $b\Delta=0$ 
\item $c\Delta0$ 
\item $\Delta=0$
\end{enumerate}
\item In the sum of first n terms of an A.P.is $cn^2$, then the sum of squares of these n terms is
\begin{enumerate}
\item $\frac{n(4n^2-1)c^2}{6}$ 
\item $\frac{n(4n^2+1)c^2}{3}$ 
\item $\frac{n(4n^2-1)c^2}{3}$ 
\item $\frac{n(4n^2+1)c^2}{6}$
\end{enumerate}
\item Let $a_1,a_2,a_3$,........be in harmonic progression with $a_1=5$ and $a_{20}=25$. The least positive integer n for which $a_n<0$ is
\begin{enumerate}
\item 22 
\item 23 
\item 24 
\item 25
\end{enumerate}
\item Let $b_1>1$ for $i=1,2,....,101.$Suppose $\log_e b_1,\log_e b_2,.....\log_e b_{101}$ are in Arithmetic progression(A.P) with the common difference $\log_e 2$. Suppose $a_1,......a_{101}$ are in A.P such that $a_1=b_1$ and $a_{51}=b_{51}$.If $t=b_1+b_2+......+b_{51}$ and $s=a_1+a_2+...a_{53},$ then 
\begin{enumerate}
\item $s>t$ and $a_{101}>b_{101}$
\item $s>t$ and $a_{101}<b_{101}$
\item $s<t$ and $a_{101}>b_{101}$
\item $s<t$ and $a_{101}<b_{101}$
\end{enumerate}

\textbf{MCQs with One or More than One Correct}

\item If the first and the $(2n-1)$st terms of an A.P., a G.P. and an H.P. are equal and their n-th terms are a,b and c respectively , then
\begin{enumerate}
\item a=b=c 
\item $a \geq b \geq c$   
\item a+c=b 
\item $ac-b^2=0$.
\end{enumerate}
\item For $0<\phi<\frac{\Pi}{2}$, if $x=\sum_{n=0}^{\infty}(\cos^{2n})\phi$, $y=\sum_{n=0}^{\infty}(\sin^{2n})\phi$, $z=\sum_{n=0}^{\infty}(\cos^{2n})\phi(\sin^{2n})\phi$ then:
\begin{enumerate}
\item xyz = xz+y 
\item xyz = xy+z 
\item xyz = x+y+z 
\item xyz = yz+x
\end{enumerate}
\item Let n be an odd integer.If $\sin n\theta=\sum_{r=0}^{n}(b)_r \sin^r\theta$, for every value of $\theta$,then 
\begin{enumerate}
\item $b_0=1, b_1=3$ 
\item $b_0=0, b_1=n$ 
\item $b_0=-1, b_1=n$ 
\item $b_0=0, b_1=n^2-3n+3$
\end{enumerate}
\item Let $T_r$ be the $r^{th}$ term of an A.P.,for r=1,2,3,..... If for some positive integers m,n we have $T_m=\frac{1}{n}$ and $T_n=\frac{1}{m}$, then $T_{mn}$ equals
\begin{enumerate}
\item $\frac{1}{mn}$ 
\item $\frac{1}{m}$+$\frac{1}{n}$ 
\item 1 
\item 0
\end{enumerate}
\item If $x>1, y>1, z>1$ are in G.P.,then $\frac{1}{1+ln x}$,$\frac{1}{1+ln y}$,$\frac{1}{1+ln z}$ are in 
\begin{enumerate}
\item A.P. 
\item H.P. 
\item G.P. 
\item None of these
\end{enumerate}
\item For a positive integer n, let $a(n)= 1+\frac{1}{2}+\frac{1}{3}+\frac{1}{4}+.....\frac{1}{(2^n)-1}$. Then
\begin{enumerate}
\item $a(100)\leq100$ 
\item $a(100)>100$ 
\item $a(200)\leq100$ 
\item $a(200)>100$
\end{enumerate}
\item A straight line through the vertex P of a triangle PQR intersects the side QR at the points 
$\vec{S}$ and the circumcircle of the triangle PQR at the point $\vec{T}$.If S is not the centre of the circumcircle, then
\begin{enumerate}
\item $\frac{1}{(PS}$+ $\frac{1}{(ST}$ $<\frac{2}{\sqrt{QS X SR}}$ 
\item $\frac{1}{(PS}$+ $\frac{1}{(ST}$ $>\frac{2}{\sqrt{QS X SR}}$  
\item $\frac{1}{(PS}$+ $\frac{1}{(ST}$ $<\frac{4}{QR}$ 
\item $\frac{1}{(PS}$+ $\frac{1}{(ST}$ $>\frac{2}{QR}$
\end{enumerate}
\item Let $S_n=\sum_{k=1}^{n} \frac{n}{n^2+kn+k^2}$ and $T_n=\sum_{k=0}^{n-1}\frac{n}{n^2+kn+k^2}$ for n=1,2,3......Then 
\item Let $S_n=\sum_{k=1}^{4n}(-1)^\frac{k(k+1)}{2}k^2$. Then $S_n$ can take value(S)
\begin{enumerate}
\item 1056 
\item 1088 
\item 1120 
\item 1332
\end{enumerate}
\item Let $\alpha$ and $\beta$ be the roots of $x^2-x-1=0$ with $\alpha > \beta$.For all positive integers n,define $a_n= \frac{\alpha^n - \beta^n}{\alpha-\beta},n\geq 1 b_1$ and $b_n=a_{n-1}+a_{n+1}, n\geq 2$ Then which of the following options is /are correct?
\begin{enumerate}
\item $\sum_{n=1}^{\infty} \frac{\alpha_n}{10^n}=\frac{10}{89}$
\item $b_n=\alpha^n+\beta^n$ for all $n\geq 1$
\item $a_1+a_2+a_3+....a_n=a_{n+2}-1$ for all $n\geq1$
\item $\sum_{n=1}^{\infty}\frac{b_n}{10^n}=\frac{8}{89}$
\end{enumerate} 

\textbf {Subjective Problems}

\item The harmonic mean of two numbers is 4. Their arithmetic mean A and the geometric mean G satisfy the relation, $2A+G^2=27$. Find the two numbers.
\item The interior angles of a polygon are in arithmetic progression. The smallest angle is 
$120\degree$, and the common difference is $5\degree$. Find the number of sides of the polygon.
\item Does there exists a geometric progression containing 27, 8 and 12 as three of its terms ? If it exits, how many such progressions are possible ?
\item Find three numbers a,b,c between 2 and 18 such that
\begin{enumerate}
\item their sum is 25 
\item the numbers 2,a,b are consecutive terms of an A.P. and  
\item the numbers b,c,18 are consecutive terms of a G.P.
\end{enumerate}
\item If $a>0,b>0$ and $c>0$, prove that $(a+b+c)(\frac{1}{a}+\frac{1}{b}+\frac{1}{c})\geq9$
\item If n is a natural number such that  $n= p_1^\alpha 1.p_2^\alpha 2.p_3^\alpha 3....p_k^\alpha k$ and $p_1,p_2,......p_k$ are distinct primes, then show that $ln n \geq k ln 2$
\item Find the sum of the series: $\sum_{r=0}^{n}(-1)^r n C_r[\frac{1}{2^r}+\frac{3^r}{2^{2r}}+\frac{7^r}{2^{3r}}+\frac{15^r}{2^{4r}}$......upto m terms]
\item Solve for x the following equation: $log_{2x+3}(6x^2+23x+21)=4-log_{3x+7}(4x^2+12x+9)$
\item If $log_3 2$, $log_3 (2^x-5)$, $log_3 (2^x-\frac{7}{2})$ are in arithmetic progression, determine the value of x.
\item Let p be the first of the n arithmetic means between two numbers and q the first of n harmonic means between the same numbers. Show that q does not lie between p and $[\frac{n+1}{n-1}]^2$p.
\item If $S_1$, $S_2$, $S_3$,......., $S_n$ are the sums of infinite geometric series whose first terms are 1, 2, 3,......, n and whose common ratios are $\frac{1}{2}$, $\frac{1}{3}$, $\frac{1}{4}$,...........$\frac{1}{n+1}$ respectively, then find the values of $S_1^2$+ $S_2^2$+ $S_3^2$,......., $S^2_{2n-1}$
\item The real numbers $x_1,x_2, x_3$ satisfying the equation $x^3-x^2+\beta x+\gamma=0$ are in A.P. Find the intervals in which $\beta$ and $\gamma$ lie.
\item Let a,b,c,d are the real numbers in G.P. If u,v,w, satisfy the system of equations $u+2v+3w=6$ $4u+5v+6w=12$ $6u+9v=4$ then show that the root of the equation $(\frac{1}{u}+\frac{1}{v}+\frac{1}{w})x^2+[(b-c)^2+(c-a)^2+(d-b)^2]x+u+v+w=0$ and $20x^2+10(a-d)^2x-9=0$ are reciprocals of the each other. 
\item The fouth power of the common difference of an arithmetic progression with integer entries is added to the product of any four consecutive terms of it. Prove that the resulting sum is the square of an integer.
\item Let $a_1$, $a_2$,.......,$a_n$ be positive real numbers in geometric progression. For each n, let $A_n$, $G_n$, $H_n$ be respectively, the arithmetic mean, geometric mean, and harmonic mean of $a_1$, $a_2$,.......,$a_n$. Find an expression for the geometric mean of $G_1$, $G_2$,......, $G_n$, in terms of $A_1$, $A_2$,.......,$A_n$,  $H_1$,$H_2$,......,$H_n$,
\item Let a, b be positive real numbers.If a, $A_1$, $A_2$, b are in arithmetic progression, a, $G_1$, $G_2$, b are in geometric progression and a, $H_1$, $H_2$, b are in harmonic progression, show that $\frac{G_1G_2}{H_1H_2}=\frac{A_1+A_2}{H_1+H_2}=\frac{(2a+b)(a+2b)}{9ab}$.
\item If a, b, c are in A.P., $a^2$, $b^2$, $c^2$ are in H.P., then prove that either a = b = c or a, b, $-\frac{c}{2}$ form a G.P.
\item If $a_n=\frac{3}{4}-[\frac{3}{4}]^2+[\frac{3}{4}]^3+......(-1)^{n-1} [\frac{3}{4}]^n$ and $b_n = 1- a_n$, then find the least natural number $n_0$ such that $b_n > a_n \forall  n \geq n_0$.

\textbf{Comprehension Based Questions}
                               
\textbf {{PASSAGE - 1}}

Let $V_r$ denote the sum of first r terms of an arithmetic progression(A.P.) whose first term is r and the common difference is (2r-1). Let $T_r=V_{r+1}-V_r-2$ and $Q_r=T_{r+1}-T_r$ for r=1,2,...

\item The sum $V_1+V_2+....+V_n$ is

\begin{enumerate}
\item $\frac{1}{12}n(n+1)(3n^2-n+1)$   
\item $\frac{1}{12}n(n+1)(3n^2+n+2)$  
\item $\frac{1}{2}n(2n^2-n+1)$   
\item $\frac{1}{3}(2n^3-2n+3)$
\end{enumerate}

\item $T_r$ is always 

\begin{enumerate}
\item an odd number          
\item an even number    
\item a prime number 
\item a composite number
\end{enumerate}

\item Which one of the following is a correct statement ?

\begin{enumerate}
\item $Q_1, Q_2, Q_3$,.... are in A.P. with common difference 5
\item $Q_1, Q_2, Q_3$,.... are in A.P. with common difference 6
\item $Q_1, Q_2, Q_3$,.... are in A.P. with common difference 11
\item $Q_1= Q_2=Q_3$=........
\end{enumerate}

\textbf {PASSAGE - 2}
               
Let $A_1, G_1, H_1$ denote the arithmetic, geometric and harmonic means respectively, of two distinct positive numbers. For $n\geq 2$, let $A_{n-1}$ and $H_{n-1}$  have  arithmetic, geometric and harmonic means as $A_n,G_n,H_n$ respectively.

\item Which one of the following statements is correct ?

\begin{enumerate}
\item $G_1 > G_2 > G_3 >..........$    
\item $G_1 < G_2 < G_3 <..........$      
\item $G_1 = G_2 = G_3 =..........$    
\item $G_1 < G_3 < G_5$ and $G_2 > G_4 > G_6 >..........$
\end{enumerate}

\item Which one of the following statements is correct ?

\begin{enumerate}
\item $A_1 > A_2 > A_3 >..........$    
\item $A_1 < A_2 < A_3 <..........$         
\item $A_1 > A_3 > A_5>....$ and $A_2 < A_4 < A_6 <..........$
\item $A_1 < A_3 < A_5<....$ and $A_2 > A_4 > A_6 >..........$
\end{enumerate}

\item Which one of the following statements is correct ?

\begin{enumerate}
\item $H_1 > H_2 > H_3 >..........$    
\item $H_1 < H_2 < H_3 <..........$        
\item $H_1 > H_3 > H_5>....$ and $H_2 < H_4 < H_6 <..........$
\item $H_1 < H_3 < H_6<....$ and $H_2 > H_4 > H_6 >..........$
\end{enumerate}
\textbf{Assertion Reson type quations} 
\item Suppose four distinct positive numbers $a_1$, $a_2$, $a_3$, $a_4$ are in G.P. Let $b_1 = a_1$, $b_2 = b_1 + a_2$, $b_3 = b_2 + a_3$ and $b_4 = b_3 + a_4$.

STATEMENT - 1: The numbers $b_1$, $b_2$, $b_3$, $b_4$ are neither in A.P. nor in G.P. and

STATEMENT - 2: The numbers $b_1$, $b_2$, $b_3$, $b_4$ are in H.P.
\begin{enumerate}
\item STATEMENT - 1 is True, STATEMENT - 2 is True; STATEMENT - 2 is a correct explanation for STATEMENT - 1
\item STATEMENT - 1 is True, STATEMENT - 2 is True; STATEMENT - 2 is a NOT a correct explanation for STATEMENT - 1
\item STATEMENT - 1 is True, STATEMENT - 2 is False
\item STATEMENT - 1 is False, STATEMENT - 2 is True
\end{enumerate}

\textbf{Integer Value Correct Type}

\item Let $S_k, k= 1, 2,......, 100$, denote the sum of the infinite geometric series whose first term is $\frac{k-1}{k!}$ and the common ratio is $\frac{1}{k}$. Then the value of $\frac{100^2}{100!}+\sum_{k=1}^{100} \abs{(k^2 - 3k + 1)S_k}$  is 
\item $a_1, a_2, a_3,.......,a_{11}$ be real numbers satisfying $a_1=15,27-2a_2>0$ and $a_k=2a_{k-1}-a_{k-2}$ for k= 3,4,.....,11, if $\frac{a_1^2 + a_2^2 +.....+a_{11}^2}{11}=90$,then the value of
$\frac{a_1 + a_2 +.....+a_{11}}{11}$ is equal to
\item Let $a_1$, $a_2$, $a_3$,.......,$a_{100}$ be an arithmetic progression with $a_1=3$ and $S_p=\sum_{i=1}^{p} a_i,1\leq p\leq 100$.For any integer n with $1\leq n\leq 20$, let m=5n .If $\frac{S_m}{S_n}$ does not depend on n, then $a_2$ is......
\item A pack contains n cards numbered from 1 to n. Two consecutive numbered cards are removed from the pack and the sum of the numbers on the remaining cards is 1224. If the smaller of the numbers on the removed cards is k, then k-20=..........
\item Let a ,b, c be positive integers such that $\frac{b}{a}$ is an integer. If a, b, c are in geometric progression and the arithmetic mean of a, b, c is b+2, then the value of 
$\frac{a^2+a-14}{a+1}$ is 
\item Suppose that all the terms of an arithmetic progression(A.P.)
are natural numbers. If the ratio of the sum of the first seven terms to the sum of the first eleven terms is 6:11 and the seventh term is lies in between 130 and 140, then the common difference of this A.P. is
\item The coefficient of $x^9$ in the expansion of $(1+x)(1+x^2)(1+x^3)......(1+x^{100})$ is
\item The sides of a right angled triangle are in arithmetic progression. If the triangle has area 24, then what is the length of its smallest side ?
\item Let X be the set consisting of the first 2018 terms of the arithmetic progression $1, 6, 11,....$, and Y be the set consisting of the first 2018 terms of the arithmetic progression $9, 16, 23,.....$ Then, the number of elements in the set $X\cup Y$ is ....
\item Let AP(a;d) denote the set of all the terms of an infinite arithmetic progression with first term a and common difference $d>0$. If AP(1;3) AP(2;5) AP(3;7)= AP(a;d) then a+d equals.....
{\textbf {Section-B JEE Main/AIEEE} }
\item If 1, $log_9(3^{1-x} +2),log_3(4.3^x - 1)$ are in A.P. then x equals
\begin{enumerate}
\item $log_3 4$  
\item $1- log_3 4$  
\item $1-log_4 3$  
\item $log_4 3$
\end{enumerate}
\item l, m, n are the $p^{th},q^{th}$ , and $r^{th}$ term of a G.P. all positive, then 
 $\myvec{\log l & p & 1 \\ \log m & q  & 1 \\ \log n n & r & 1}$  equals
\begin{enumerate}
\item -1  
\item 2   
\item 1  
\item 0
\end{enumerate}
\item The value of $2^\frac{1}{4}.4^\frac{1}{8}.8^\frac{1}{16}.........\infty$ is
\begin{enumerate}
\item 1  
\item 2   
\item $\frac{3}{2}$  
\item 4
\end{enumerate}
\item Fifth term of a G.P. is 2, then the product of its 9 terms is
\begin{enumerate}
\item 256  
\item 512   
\item 1024  
\item none of these
\end{enumerate}
\item Sum of infinite number of terms of GP is 20 and sum of their square is 100. The common ratio of GP is
\begin{enumerate}
\item 5  
\item $\frac{3}{5}$   
\item $\frac{8}{5}$  
\item $\frac{1}{5}$
\end{enumerate}
\item $1^3-2^3+3^3-4^3+....+9^3$=
\begin{enumerate}
\item 425  
\item -425   
\item 475  
\item -475
\end{enumerate}
\item The sum of the series $\frac{1}{1.2}-\frac{1}{2.3}+\frac{1}{3.4}.......$ upto $\infty$ is equal to 
\begin{enumerate}
\item $log_e\frac{4}{e}$
\item $2 log_e 2$   
\item $log_e 2-1$  
\item $log_e 2$
\end{enumerate}
\item If $S_n = \sum_{r=0}^{n}\frac{1}{nC_r}$ and $t_n = \sum_{r=0}^{n}\frac{r}{nC_r}$, then 
$\frac{t_n}{S_n}$ is equal to 
\begin{enumerate}
\item $\frac{2n-1}{2}$
\item $\frac{1}{2}$n-1   
\item  n-1 
\item $\frac{1}{2}$n
\end{enumerate}
\item Let $T_r$ be the $r^{th}$ term of an A.P.whose first term is a and common difference is d. If for some positive integers m, n, $m\neq n$, $T_m = \frac{1}{n}$ and $T_n = \frac{1}{m}$, then a-d equals
\begin{enumerate}
\item $\frac{1}{m}+\frac{1}{n}$
\item 1
\item $\frac{1}{mn}$
\item 0
\end{enumerate}
\item The sum of the first n terms of the series $1^2 + 2.2^2 + 3^2 + 2.4^2 + 5^2 + 2.6^2+.....$ is 
$\frac{n(n+1)^2}{2}$ when n is even. When n is odd the sum is 
\begin{enumerate}
\item $[\frac{n(n+1)}{2}]^2$
\item $\frac{n^2(n+1)}{2}$
\item $\frac{n(n+1)^2}{4}$
\item $\frac{3n(n+1)}{2}$
\end{enumerate}
\item The sum of series $\frac{1}{2!}+\frac{1}{4!}+\frac{1}{6!}.....$ is 
\begin{enumerate}
\item $\frac{(e^2 - 2)}{e}$
\item $\frac{(e-1)^2}{2e}$
\item $\frac{(e^2-1)}{2e}$
\item $\frac{(e^2-1)}{2}$
\end{enumerate}
\item If the coefficient of $r^{th}$,$(r+1)^th$ and $(r+2)^{th}$ terms in the binomial expansion of $(1+y)^m$ are in A.P., then m and r satisfy the equation
\begin{enumerate}
\item $m^2 - m(4r-1) + 4r^2 - 2 = 0$
\item $m^2 - m(4r+1) + 4r^2 + 2 = 0$
\item $m^2 - m(4r+1) + 4r^2 - 2 = 0$
\item $m^2 - m(4r-1) + 4r^2 + 2 = 0$
\end{enumerate}
\item If $x=\sum_{n=0}^{\infty}a^n$,  $y=\sum_{n=0}^{\infty}b^n$,  $z=\sum_{n=0}^{\infty}c^n$ where a, b, c are in A.P. and $\abs{a}<1,\abs{b}<1,\abs{c}<1$ then x, y, z are in
\begin{enumerate}
\item G.P.
\item A.P.
\item Arithmetic - Geometric Progression
\item H.P.
\end{enumerate}
\item The sum of series 1+$\frac{1}{4.2!}+\frac{1}{16.4!}+\frac{1}{64.6!}.....\infty$ is 
\begin{enumerate}
\item $\frac{(e - 1)}{\sqrt{e}}$
\item $\frac{(e + 1)}{\sqrt{e}}$
\item $\frac{(e - 1)}{2\sqrt{e}}$
\item $\frac{(e + 1)}{2\sqrt{e}}$
\end{enumerate}
\item Let $a_1$, $a_2$, $a_3$.....be terms in A.P. If $\frac{a_1+a_2+......+a_p}{a_1+a_2+......+a_q} = \frac{p^2}{q^2}$, $p\neq q$, then $\frac{a_6}{a_{21}}$ equals 
\begin{enumerate}
\item $\frac{41}{11}$
\item $\frac{7}{2}$.
\item $\frac{2}{7}$
\item $\frac{11}{41}$
\end{enumerate}
\item If $a_1$, $a_2$, $a_3$.....$a_n$ are in H.P., then the expression $a_1a_2+a_2a_3+.....+a_{n-1}a_n$ is equal to
\begin{enumerate}
\item $n(a_1 -a_n)$
\item $(n-1)(a_1 -a_n)$
\item $n(a_1a_n)$
\item $(n-1)(a_1a_n)$
\end{enumerate}
\item The sum of series $\frac{1}{2!}-\frac{1}{3!}+\frac{1}{4!}.....$ upto $\infty$ is 
\begin{enumerate}
\item $e^-\frac{1}{2}$
\item $e^+\frac{1}{2}$
\item $e^{-2}$
\item $e^{-1}$
\end{enumerate}
\item In a geometric progression consisting of positive terms, each term equals the sum of the next two terms. Then the common ratio of its progression is equals
\begin{enumerate}
\item $\sqrt{5}$
\item $\frac{1}{2}(\sqrt{5} - 1)$
\item $\frac{1}{2}(1-\sqrt{5})$
\item $\frac{1}{2}\sqrt{5}$
\end{enumerate}
\item The first two terms of a geometric progression add up to 12. the sum of the third and the fourth terms is 48. If the terms of the geometric progression are alternately positive and negative, then the first term is
\begin{enumerate}
\item -4
\item -12
\item 12
\item 4
\end{enumerate}
\item The sum to infinite term of the series $1+\frac{2}{3}+\frac{6}{3^2}+\frac{10}{3^3}+\frac{14}{3^4}+.....$ is 
\begin{enumerate}
\item 3
\item 4
\item 6
\item 2
\end{enumerate}
\item A person is to count 4500 currency notes. Let $a_n$ denote the number of notes he counts in the $n^{th}$ minute. If $a_1=a_2=.....=a_{10}=150$ and $a_{10},a_{11},.....$ are in A.P. with common difference -2, then the time taken by him to count all notes is
\begin{enumerate}
\item 34 minutes
\item 125 minutes
\item 135 minutes
\item 24 minutes
\end{enumerate}
\item A man saves 200 in each of the first three months of his service. In each of the subsequent months his saving increases by 40 more than the saving of immediately previous month. His total savings from the start of service will be 11040 after
\begin{enumerate}
\item 19 months
\item 20 months
\item 21 months
\item 18 months
\end{enumerate}
\item \textbf{Statement - 1:} The sum of the series $1+(1+2+4)+(4+6+9)+(9+12+16)+.....+(361+380+400)$ is $8000.$\\
	  \textbf{Statement - 2:} $\sum_{k=1}^{n}(k^3 -(k-1)^3)=n^3$, for any natural number n.

\begin{enumerate}
\item Statement - 1 is false, Statement - 2 is true.
\item Statement - 1 is true, Statement - 2 is true, Statement - 2 is a correct explanation for Statement - 1
\item Statement - 1 is true, Statement - 2 is true, Statement - 2 is a not a correct explanation for Statement - 1
\item Statement - 1 is true, Statement - 2 is false.
\end{enumerate}
\item The sum of the first 20 terms of sequence 0.7,0.77,0.777,....,is
\begin{enumerate}
\item $\frac{7}{81}(179-10^{-20})$
\item $\frac{7}{9}(99-10^{-20})$
\item $\frac{7}{81}(179+10^{-20})$
\item $\frac{7}{9}(99+10^{-20})$
\end{enumerate}
\item If $(10)^9+2(11)^1(10^8)+3(11)^2(10)^7+........+10(11)^9=k(10)^9,$ then k is equal to :
\begin{enumerate}
\item 100
\item 110
\item $\frac{122}{10}$
\item $\frac{441}{100}$
\end{enumerate}
\item Three positive numbers form an increasing G.P. If the middle term in this G.P. is doubled, the new numbers are in A.P. then the common ration of the G.P. is:
\begin{enumerate}
\item $2-\sqrt{3}$
\item $2+\sqrt{3}$
\item $\sqrt{2}+\sqrt{3}$
\item $3+\sqrt{2}$
\end{enumerate}
\item The sum of the first 9 terms of the series.$\frac{1^3}{1}+\frac{1^3+2^3}{1+3}+\frac{1^3+2^3+3^3}{1+3+5}$+.......
\begin{enumerate}
\item 142
\item 192
\item 71
\item 96
\end{enumerate}
\item If m is the A.M. of two distinct real numbers l and $n(l,n>1)$ and $G_1,G_2$ and $G_3$ are the three geometric means between l and n, then $G_1^4+2G_2^4+G_3^4$ equals:
\begin{enumerate}
\item $4lmn^2$
\item $4l^2m^2n^2$
\item $4l^2mn$
\item $4lm^2n$
\end{enumerate}
\item If the $2^{nd},5^{th}$ and $9^{th}$ terms of a non-constant A.P. are in G.P., then the common ratio of this G.P. is:
\begin{enumerate}
\item 1
\item $\frac{7}{4}$
\item $\frac{8}{5}$
\item $\frac{4}{3}$
\end{enumerate}
\item If the sum of the first ten terms of the series $(1\frac{3}{5})^2+(2\frac{2}{5})^2+(3\frac{1}{5})^2+4^2+(4\frac{4}{5})^2$+......, is $\frac{16}{5}$m then m is equal to :
\begin{enumerate}
\item 100
\item 99
\item 102
\item 101
\end{enumerate}
\item If, for a positive integer n, the quadratic equation, $x(x+1)+(x+1)(x+2)+.....+(x+\overline{n-1})(x+n)=10n$ has two consecutive integral solutions, then n is equal to :
\begin{enumerate}
\item 11
\item 12
\item 9
\item 10
\end{enumerate}
\item For any three positive real numbers a, b and c, 
\begin{align} 
9(25a^2+b^2)+25(c^2-3ac)=15b(3a+c)
\end{align} Then:
\begin{enumerate}
\item a, b and c are in G.P.
\item b, c and a are in G.P.
\item b, c and a are in A.P.
\item a, b and c are in A.P.
\end{enumerate}
\item Let a, b, c $\in$ R. If $f(x)=ax^2+bx+c$ is such that $a+b+c=3$   $f(x+y)=f(x)+f(y)+xy,\forall x, y \in R$, then $\sum_{n=1}^{10}f(n)$ is equal to :
\begin{enumerate}
\item 255
\item 330
\item 165
\item 190
\end{enumerate}
\item Let $a_1,a_2,a_3,......,a_{49}$ be in A.P. such that $\sum_{k=0}^{12}a_{4k+1}=416$ and $a_9+a_{43}=66$.If $a_1^2+a_2^2+.....+a_{17}^2=140m$, then m is equal to
\begin{enumerate}
\item 68
\item 34
\item 33
\item 66
\end{enumerate}
\item Let A be the sum of the first 20 terms and B be the sum of the first 40 terms of the series $1^2+2.2^2+3^2+2.4^2+5^2+2.6^2+.....$ If $B-2A=100\Lambda$, then $\Lambda$ is equal to :
\begin{enumerate}
\item 248
\item 464
\item 496
\item 232
\end{enumerate}
\item If a, b, and c be three distinct real numbers in G.P. and a+b+c=xb, then x cannot be :
\begin{enumerate}
\item -2
\item -3
\item 4
\item 2
\end{enumerate}
\item Let $a_1,a_2,a_3,......,a_{30}$ be in A.P., $S=\sum_{i=1}^{30}a_i$ and $T=\sum_{i=1}^{15}a_{2i-1}$. If $a_5=27$ and S-2T=75, Then $a_{10}$ is equal to :
\begin{enumerate}
\item 52
\item 57
\item 47
\item 42
\end{enumerate}
\item Three circles of radii a, b, c $(a < b < c)$ touch each other externally. If they have X-axis as a common tangent, then :
\begin{enumerate}
\item $\frac{1}{\sqrt{a}}=\frac{1}{\sqrt{b}}+\frac{1}{\sqrt{c}}$
\item $\frac{1}{\sqrt{b}}=\frac{1}{\sqrt{a}}+\frac{1}{\sqrt{c}}$
\item a, b, c are in A.P.
\item $\sqrt{a}, \sqrt{b} ,\sqrt{c}$ are in A.P.
\end{enumerate}
\item Let the sum of the first n terms of a non-constant A.P., $a_1,a_2,a_3,..........$ be 
$50n+\frac{n(n-7)}{2}A$, where A is a constant. If d is the common difference of this A.P.,then the ordered pair $(d,a_{10})$ is equal to:
\begin{enumerate}
\item (50, 50+46A)
\item (50, 50+45A)
\item (A, 50+45A)
\item (A, 50+46A)
\end{enumerate}
\end{enumerate} 
 

\section{JEE Exercises:Mathematical Induction and Binomial Theorem}
%\documentclass[journal,12pt,twocolumn]{IEEEtran}
%\usepackage{setspace}
%\usepackage{gensymb}
%\usepackage{caption}
%\usepackage{subfiles}
%%\usepackage{multirow}
%%\usepackage{multicolumn}
%%\usepackage{subcaption}
%%\doublespacing
%\singlespacing
%\usepackage{csvsimple}
%\usepackage{amsmath}
%\usepackage{multicol}
%%\usepackage{enumerate}
%\usepackage{amssymb}
%%\usepackage{graphicx}
%\usepackage{newfloat}
%%\usepackage{syntax}
%\usepackage{listings}
%%\usepackage{iithtlc}
%\usepackage{color}
%\usepackage{tikz}
%\usetikzlibrary{shapes,arrows}
%
%
%
%%\usepackage{graphicx}
%%\usepackage{amssymb}
%%\usepackage{relsize}
%%\usepackage[cmex10]{amsmath}
%%\usepackage{mathtools}
%%\usepackage{amsthm}
%%\interdisplaylinepenalty=2500
%%\savesymbol{iint}
%%\usepackage{txfonts}
%%\restoresymbol{TXF}{iint}
%%\usepackage{wasysym}
%\usepackage{amsthm}
%\usepackage{mathrsfs}
%\usepackage{txfonts}
%\usepackage{stfloats}
%\usepackage{cite}
%\usepackage{cases}
%\usepackage{mathtools}
%\usepackage{caption}
%\usepackage{enumerate}
%\usepackage{tfrupee}	
%\usepackage{enumitem}
%\usepackage{amsmath}
%%\usepackage{xtab}
%\usepackage{longtable}
%\usepackage{multirow}
%%\usepackage{algorithm}
%%\usepackage{algpseudocode}
%\usepackage{enumitem}
%\usepackage{mathtools}
%\usepackage{hyperref}
%%\usepackage[framemethod=tikz]{mdframed}
%\usepackage{listings}
%    %\usepackage[latin1]{inputenc}                                 %%
%    \usepackage{color}                                            %%
%    \usepackage{array}                                            %%
%    \usepackage{longtable}                                        %%
%    \usepackage{calc}                                             %%
%    \usepackage{multirow}                                         %%
%    \usepackage{hhline}                                           %%
%    \usepackage{ifthen}                                           %%
%  %optionally (for landscape tables embedded in another document): %%
%    \usepackage{lscape}     
%
%
%\usepackage{url}
%\def\UrlBreaks{\do\/\do-}
%
%
%%\usepackage{stmaryrd}
%
%
%%\usepackage{wasysym}
%%\newcounter{MYtempeqncnt}
%\DeclareMathOperator*{\Res}{Res}
%%\renewcommand{\baselinestretch}{2}
%\renewcommand\thesection{\arabic{section}}
%\renewcommand\thesubsection{\thesection.\arabic{subsection}}
%\renewcommand\thesubsubsection{\thesubsection.\arabic{subsubsection}}
%
%\renewcommand\thesectiondis{\arabic{section}}
%\renewcommand\thesubsectiondis{\thesectiondis.\arabic{subsection}}
%\renewcommand\thesubsubsectiondis{\thesubsectiondis.\arabic{subsubsection}}
%
%% correct bad hyphenation here
%\hyphenation{op-tical net-works semi-conduc-tor}
%
%%\lstset{
%%language=C,
%%frame=single, 
%%breaklines=true
%%}
%
%%\lstset{
%	%%basicstyle=\small\ttfamily\bfseries,
%	%%numberstyle=\small\ttfamily,
%	%language=Octave,
%	%backgroundcolor=\color{white},
%	%%frame=single,
%	%%keywordstyle=\bfseries,
%	%%breaklines=true,
%	%%showstringspaces=false,
%	%%xleftmargin=-10mm,
%	%%aboveskip=-1mm,
%	%%belowskip=0mm
%%}
%
%%\surroundwithmdframed[width=\columnwidth]{lstlisting}
%\def\inputGnumericTable{}                                 %%
%\lstset{
%%language=C,
%frame=single, 
%breaklines=true,
%columns=fullflexible
%}
% 
%
%\begin{document}
%%
%\tikzstyle{block} = [rectangle, draw,
%    text width=3em, text centered, minimum height=3em]
%\tikzstyle{sum} = [draw, circle, node distance=3cm]
%\tikzstyle{input} = [coordinate]
%\tikzstyle{output} = [coordinate]
%\tikzstyle{pinstyle} = [pin edge={to-,thin,black}]
%
%\theoremstyle{definition}
%\newtheorem{theorem}{Theorem}[section]
%\newtheorem{problem}{Problem}
%\newtheorem{proposition}{Proposition}[section]
%\newtheorem{lemma}{Lemma}[section]
%\newtheorem{corollary}[theorem]{Corollary}
%\newtheorem{example}{Example}[section]
%\newtheorem{definition}{Definition}[section]
%%\newtheorem{algorithm}{Algorithm}[section]
%%\newtheorem{cor}{Corollary}
%\newcommand{\BEQA}{\begin{eqnarray}}
%\newcommand{\EEQA}{\end{eqnarray}}
%\newcommand{\define}{\stackrel{\triangle}{=}}
%
%\bibliographystyle{IEEEtran}
%%\bibliographystyle{ieeetr}
%
%\providecommand{\nCr}[2]{\,^{#1}C_{#2}} % nCr
%\providecommand{\nPr}[2]{\,^{#1}P_{#2}} % nPr
%\providecommand{\mbf}{\mathbf}
%\providecommand{\pr}[1]{\ensuremath{\Pr\left(#1\right)}}
%\providecommand{\qfunc}[1]{\ensuremath{Q\left(#1\right)}}
%\providecommand{\sbrak}[1]{\ensuremath{{}\left[#1\right]}}
%\providecommand{\lsbrak}[1]{\ensuremath{{}\left[#1\right.}}
%\providecommand{\rsbrak}[1]{\ensuremath{{}\left.#1\right]}}
%\providecommand{\brak}[1]{\ensuremath{\left(#1\right)}}
%\providecommand{\lbrak}[1]{\ensuremath{\left(#1\right.}}
%\providecommand{\rbrak}[1]{\ensuremath{\left.#1\right)}}
%\providecommand{\cbrak}[1]{\ensuremath{\left\{#1\right\}}}
%\providecommand{\lcbrak}[1]{\ensuremath{\left\{#1\right.}}
%\providecommand{\rcbrak}[1]{\ensuremath{\left.#1\right\}}}
%\theoremstyle{remark}
%\newtheorem{rem}{Remark}
%\newcommand{\sgn}{\mathop{\mathrm{sgn}}}
%\providecommand{\abs}[1]{\left\vert#1\right\vert}
%\providecommand{\res}[1]{\Res\displaylimits_{#1}} 
%\providecommand{\norm}[1]{\left\Vert#1\right\Vert}
%\providecommand{\mtx}[1]{\mathbf{#1}}
%\providecommand{\mean}[1]{E\left[ #1 \right]}
%\providecommand{\fourier}{\overset{\mathcal{F}}{ \rightleftharpoons}}
%%\providecommand{\hilbert}{\overset{\mathcal{H}}{ \rightleftharpoons}}
%\providecommand{\system}{\overset{\mathcal{H}}{ \longleftrightarrow}}
%	%\newcommand{\solution}[2]{\textbf{Solution:}{#1}}
%\newcommand{\solution}{\noindent \textbf{Solution: }}
%\newcommand{\myvec}[1]{\ensuremath{\begin{pmatrix}#1\end{pmatrix}}}
%\providecommand{\dec}[2]{\ensuremath{\overset{#1}{\underset{#2}{\gtrless}}}}
%\DeclarePairedDelimiter{\ceil}{\lceil}{\rceil}
%%\numberwithin{equation}{section}
%%\numberwithin{problem}{subsection}
%%\numberwithin{definition}{subsection}
%\makeatletter
%\@addtoreset{figure}{section}
%\makeatother
%
%\let\StandardTheFigure\thefigure
%%\renewcommand{\thefigure}{\theproblem.\arabic{figure}}
%\renewcommand{\thefigure}{\thesection}
%
%
%%\numberwithin{figure}{subsection}
%
%%\numberwithin{equation}{subsection}
%%\numberwithin{equation}{section}
%%\numberwithin{equation}{problem}
%%\numberwithin{problem}{subsection}
%\numberwithin{problem}{section}
%%%\numberwithin{definition}{subsection}
%%\makeatletter
%%\@addtoreset{figure}{problem}
%%\makeatother
%\makeatletter
%\@addtoreset{table}{section}
%\makeatother
%
%\let\StandardTheFigure\thefigure
%\let\StandardTheTable\thetable
%\let\vec\mathbf
%%%\renewcommand{\thefigure}{\theproblem.\arabic{figure}}
%%\renewcommand{\thefigure}{\theproblem}
%
%%%\numberwithin{figure}{section}
%
%%%\numberwithin{figure}{subsection}
%
%
%
%\def\putbox#1#2#3{\makebox[0in][l]{\makebox[#1][l]{}\raisebox{\baselineskip}[0in][0in]{\raisebox{#2}[0in][0in]{#3}}}}
%     \def\rightbox#1{\makebox[0in][r]{#1}}
%     \def\centbox#1{\makebox[0in]{#1}}
%     \def\topbox#1{\raisebox{-\baselineskip}[0in][0in]{#1}}
%     \def\midbox#1{\raisebox{-0.5\baselineskip}[0in][0in]{#1}}
%
%\vspace{3cm}
%
%\title{ 
%%	\logo{
%Mathematical Induction and Binomial Theorem
%%	}
%}
%
%\author{ G V V Sharma$^{*}$% <-this % stops a space
%	\thanks{*The author is with the Department
%		of Electrical Engineering, Indian Institute of Technology, Hyderabad
%		502285 India e-mail:  gadepall@iith.ac.in. All content in this manual is released under GNU GPL.  Free and open source.}
%	
%}	
%
%\maketitle
%
%%\tableofcontents
%
%\bigskip
%
%\renewcommand{\thefigure}{\theenumi}
%\renewcommand{\thetable}{\theenumi}
%
%
%
%\begin{enumerate}[label=\arabic*]
%\numberwithin{equation}{enumi}
\renewcommand{\theequation}{\theenumi}
\begin{enumerate}[label=\arabic*.,ref=\thesubsection.\theenumi]
\numberwithin{equation}{enumi}

\item The large of $99^{50}+100^{50}$ and $101^{50}$ is........
\item The sum of the coefficients of the polynomial$(1+x-3x^2)^{2163}$ is....
\item If
\begin{align}
(1+ax)^n=1+8x+24x^2+.....
\end{align} then a=....and n=.....\\
\item let n be positive integer. If the coefficients of 2nd,3rd and 4th terms in the expansion of $(1+x)^n$ are in A.P.,then value of n is.....
\item the sum of the rational terms in the expansion of  $(\sqrt{2}+3^\frac{1}{5})^{10}$ is .......
\item Given Positive integers r$>1$,n$>1$ and that the coefficient of (3r)th and (r+2)th terms in the binomial expression of $(1+x)^{2n}$ are equal.Then
\begin{enumerate}
\item n=2r

\item n=3r

\item n=2r+1

\item none of these
\end{enumerate}

\item The Coefficient of $x^4$ in$ (\frac{x}{2}-\frac{3}{x^2})^{10} $ is
\begin{enumerate}
\item $\frac{405}{256}$

\item $\frac{504}{259}$

\item $\frac{450}{263}$

\item none of these
\end{enumerate}

\item The expression $(x+(x^3-1)^\frac{1}{2})^5+(x-(x^3-1)^\frac{1}{2})^5$ is a polynomial of degree 
\begin{enumerate}
\item 5

\item 6

\item 7

\item 8
\end{enumerate}
\item If in the expression of $(1+x)^m (1-x)^n$,the coefficient of x and $x^2$ are 3 and -6 respectively,then m is
\begin{enumerate}
\item 6

\item 9

\item 12

\item 24
\end{enumerate}
\item For $2 \leq r \leq n$,
$\myvec{n\\r}+2\myvec{n\\r-1}+\myvec{n\\r-2}=$
\begin{enumerate}
 \item $\myvec{n+1\\r-1}$

 \item $2\myvec{n+1\\r+1}$

 \item $2\myvec{n+2\\r}$

 \item $\myvec{n+2\\r}$ 
 \end{enumerate}
\item In the binomial expression of $(a-b)^n$,n$\geq$5,the sum of the 5th and 6th terms is Zero.Then $\frac{a}{b}$ equals
\begin{enumerate}
\item $\frac{n-5}{6}$

\item $\frac{n-4}{5}$

\item $\frac{5}{n-4}$

\item $\frac{6}{n-5}$
\end{enumerate}
\item The sum $$\sum_{i=0}^{m} \myvec{10\\i} \myvec{20\\m-i}, $$ (where \myvec{p\\q}=0 if p<q) is maximum when m is
\begin{enumerate}
\item 5

\item 10

\item 15

\item 20
\end{enumerate}
\item Coefficient of $t^{24}$ in $(1+t^2)^{12}(1+t^{12})(1+t^{24})$ is
\begin{enumerate}
\item $^{12}C_6+3$

\item $^{12}C_6+1$

\item $^{12}C_6$

\item $^{12}C_6+2$
\end{enumerate}
\item If $^{n-1}C_r=(k^2-3) ^{n}C_{r+1}$,Then k$\in$
\begin{enumerate}
\item $(-\infty,-2]$
\item $(2,-\infty,-2]$
\item $[-\sqrt{3},\sqrt{3}]$
\item $(\sqrt{3},-2]$
\end{enumerate}
\item The value of $\myvec{30\\0}\myvec{30\\10}-\myvec{30\\1}\myvec{30\\11}+\myvec{30\\2}\myvec{30\\12}....+\myvec{30\\20}\myvec{30\\30}$ is where $\myvec{n\\r}=^nC_r$
\begin{enumerate}
\item $\myvec{30\\10}$

\item $\myvec{30\\15}$

\item $\myvec{60\\30}$

\item $\myvec{31\\10}$
\end{enumerate}
 \item For r=0,1,....10,let $A_r,B_r and C_r$ denote respectively,the coefficient of $x^r$ in the expansions of $(1+x)^{10},(1+x)^{20} and (1+x)^{30}$.Then $$\sum_{r=1}^{10} A_r(B_{10}B_r-C_{10}A_r)$$ is equal to
 \begin{enumerate}
 \item $(B_{10}-C_{10})$
 \item $A_{10}(B^2_{10} C_{10} A_{10})$
 \item 0
 \item $(C_{10}-B_{10})$
 \end{enumerate}
 \item Coefficient of $x^{11}$ in the expansion of $(1+x^2)^4$ $(1+x^3)^7$ $(1+x^4)^{12}$ is
 \begin{enumerate}
     \item 1051
     \item 1106
     \item 1113
     \item 1120
 \end{enumerate}
\item If $C_r$ stands for $^{n}C_r$,then the sum of the series $\frac{2(\frac{n}{2})!(\frac{n}{2})!}{n!}[C^2_0-2C^2_1+3C^2_2-.....+(-1)^n(n+1)C^2_n]$,where n is an even positive integer,is equal to 
\begin{enumerate}
    \item 0
    \item $(-1)^\frac{n}{2}(n+1)$
    \item $(-1)^\frac{n}{2}(n+2)$
    \item $(-1)^n n$
    \item none of these
\end{enumerate}
\item if $a_n=\sum_{r=0}^{n} \frac{1}{^{n}C_r}$,then $\sum_{r=0}^{n} \frac{r}{^{n}C_r}$ equals
\begin{enumerate}
    \item $(n-1)a_n$
    \item $na_n$
    \item $\frac{1}{2}na_n$
    \item none of these 
\end{enumerate}
\item Given that 
\begin{align}
C_1+2C_2x+3C_3x^2+......+2nC_{2n}x^{2n-1}=2n(1+x)^{2n-1}
\end{align} where $C_r=\frac{(2n!)}{r!(2n-r)!}$  r=0,1,2,....,2n Prove that $C^2_1-2C^2_2+3C^2_3-......-2nC^2_2n=(-1)^{n}nC_n$.
\item Prove that $7^{2n}+(2^{3n-3})(3^{n-1})$ is divisible by 25 for any natural numbers n
\item If 
\begin{align}
(1+x)^n=C_0+C_1x+C_2x^2+......+C_nx^n
\end{align} then show that the sum of the products of the $C_i$'s taken at a time,represented by $\sum_{0\leq i<j \leq n}\sum C_i C_j$ is equal to $2^{2n-1}-\frac{(2n)!}{2(n!)^2}$
\item Use mathematical induction to prove : If n is any odd positive integer,then $n(n^2-1)$ is divisible by 24.
\item If p be a natural number then prove that $p^{n+1}+(p+1)^{2n-1}$ is divisible by $p^2+p+1$ for every position integer n.
\item Given $s_n=1+q+q^2+....+q^n$;
$s_n=1+\frac{q+1}{2}+(\frac{q+1}{2})^2+......+(\frac{q+1}{2})^n$, $q\neq 1$
prove that $^{n+1}C_1+^{n+1}C_2+ ^{n+1}C_3+.....+^{n+1}C_n s_n=2^n S_n$
\item Use method of mathematical induction that $2.7^n+3.5^n-5$ is divisible by 24 for all n$>$0
\item Prove by mathematical induction that $\frac{(2n)!}{2^{2n}(n!)^2}$ $\leq$ $\frac{1}{(3n+1)^\frac{1}{2}}$ for all positive integers n.
\item Let R=$(5\sqrt{5}+11)^{2n+1}$ and f=R-[R],where [] denotes the greater integer function.Prove that $Rf=4^{2n+4}$
\item Using mathematical induction prove that $^{m} C_0  ^{n} C_k+ ^{m} C_1  ^{n} C_{k-1}+.....+ ^{m} C_k  ^{n}C_0= ^{m+n} C_k$,where m,n,k are positive integers,and $^{p}C_q=0$ for p$<$q.
\item Prove that $C_0-2^2C_1+3^2C_2-......+(-1)^n (n+1)^2 C_n=0$,n$>$2,where $C_r=^{n}C_r$.
\item Prove that $\frac{n^7}{7}+\frac{n^5}{5}+\frac{2n^3}{3}-\frac{n}{105}$ is a integer for every positive integer n.
\item using induction or otherwise ,Prove that for any non-negative integers m,n,r and k,
$\sum_{m=0}^{k} (n-m)\frac{(r+m)!}{m!}=\frac{(r+k+1)!}{k!}[\frac{n}{r+1}-\frac{k}{r+2}]$
\item If 
\begin{align}
\sum_{r=0}^{2n}a_r(x-2)^r=\sum_{r=0}^{2n}b_r(x-3)^r
\end{align} and $a_k=1$ for all k $\geq$ n,then show that $b_n= ^{2n+1}C_{n+1}$
\item Let p$\geq$ 3 be an Integer and $\alpha,\beta$ be the roots of 
\begin{align}x^2-(p+1)x+1=0
\end{align} Using mathematical induction show that $\alpha^n+\beta^n$
(i) is an integer and (ii) is not divisible by p
\item Using mathematical induction Prove that \\
$\tan^{-1}\frac{1}{3}+\tan^{-1}\frac{1}{7}+....\tan^{-1}{1/(n^2+n+1)}=\tan^{-1}{n/(n+2)}$
\item Prove that $\sum_{r=1}^{k}(-3)^{r-1}\enspace {^{3n}C_{2r-1}}=0$, where $k=\frac{3n}{2}$ and n is an even positive integer.
\item If x is not an integral multiple of 2$\pi$ use mathematical induction to prove that:
$\cos x+\cos 2x+....+\cos nx=\cos \frac{n+1}{2}x\sin \frac{nx}{2}$ cosec$\frac{x}{2}$
\item let n be positive integer and \begin{align} (1+x+x^2)=a_0+a_1x+.....+a_{2n}x^{2n}\end{align} show that $a^2_0-a^1_2+a^2_2.....+a^2_{2n}=a_n$
\item Using mathematical induction prove that for every integer n$\geq$ 1,$(3^{2n}-1)$ is divisible by $2^{n+2}$ but not by $2^{n+3}$
\item Let $0<A_i<\pi$ for i=1,2,....,n.Use mathematical induction to prove that $\sin A_1+\sin A_2...+\sin A_n \leq n \sin(\frac{A_1+A_2+...+A_n}{n})$ where $\geq$ 1 is natural number.\\
(You may use the fact that $p\sin x+(1-p) \sin y\leq \sin [px+(1-p)y]$,where 0$\leq $p$\leq$1 and 0$\leq $x,y$\leq\pi$)
\item Let p be prime and m a positive integer.By mathematical induction on m or otherwise prove that whenever r is an integer such that p does not divide r,p divides $^{mp}C_r$,
[Hint: you may use the fact that $(1+x)^{(m+1)p}=(1+x)^p (1+x)^{mp}$]
\item Let n be any positive integer.prove that $\sum_{k=0}^{m} \frac{\myvec{2n-k\\k}}{\myvec{2n-k\\n}}.\frac{(2n-4k+1)}{(2n-2k+1)}2^{n-2k}=\frac{\myvec{n\\m}}{\myvec{2n-2m\\n-m}}2^{n-2m}$ for each non-be gatuve integer m $\leq $ n.\\(Here $\myvec{p\\q}=^{p}C_q$)
\item For any Positive integer m, n(with n $\geq$ m),let $\myvec{n\\m}= ^{n}C_m$. Prove that $\myvec{n\\m}+\myvec{n-1\\m}+\myvec{n-2\\m}+....+\myvec{m\\m}=\myvec{n+1\\m+2}$ Hence or otherwise, prove that$\myvec{n\\m}+2\myvec{n-1\\m}+3\myvec{n-2\\m}+....+(n-m+1)\myvec{m\\m}=\myvec{n+2\\m+2}$
\item for every positive integer n,Prove that\\
 $\sqrt{(4n+1)}<\sqrt{n}+\sqrt{n+1}<\sqrt{4n+2}.$ Hence or otherwise, prove that\\
  $[\sqrt{n}+\sqrt{(n+1)}]=[\sqrt{4n+1}]$ where [x] denotes the greater integers not exceeding x.
\item Let a,b,c be the positive real numbers such that $b^2-4ac>0$ and let $\alpha_1$=c prove by induction that$\alpha{n+1} = \frac{a\alpha^2_n}{(b^2-2a(\alpha_1+\alpha_2+...+\alpha_n))}$is well defined and 
$\alpha{n+1} < \frac{\alpha_n}{2}$ for all n = 1,2,.....(Here, 'well-defined' means that the denominator in the expresion for $\alpha_{n+1}$ is not zero.) 

\item Use the mathematical induction to show that $(25)^{n+1}-24n+5735$ is divisible by $(24)^2$ for all n=1,2,......
\item prove that $2^k\myvec{n\\0}\myvec{n\\k}-2^{k-1\myvec{n\\2}}\myvec{n\\1}\myvec{n-1\\k-1}+2^{k-2}\myvec{n-2\\k-2}-....(-1)^k \myvec{n\\k}\myvec{n-k\\0}=\myvec{n\\k}$
\item A coin has probability p of showing head when tossed.It is tossed n times.Let $p_n$ denote the probability that no two (or more) consecutive heads occur Prove that $p_1=1,p_2=1-p^2$ and $p_n=(1-p).p_(n-1)+p(1-p)p_{n-2}$ for all n $\geq$ 3.\\
Prove by induction on n, that $p_n=A\alpha^n+B\beta^n$ for all n $\geq$ 1,\\
where $\alpha$ and $\beta$ are the roots of quadratic equation 
\begin{align}
x^2-(1-p)x-p(1-p)=0
\end{align} and $A=\frac{p^2+\beta-1}{\alpha\beta-\alpha^2},B=\frac{p^2+\alpha-1}{\alpha\beta-\beta^2}.$

\item The coefficient of three consecutive terms of $(1+x)^{n+5}$ are in the ratio 5:10:14.Then n=
\item Let m be the smallest positive integer such that the Coefficient of $x^2$ in the expansion of \\$(1+x)^2+(1+x)^3+.....+(1+x)^{49}+(1+mx)^{50}$ is $(3n+1)^{51}C_3$ for some positive integers n.Then value of n is
\item Let X=$(^{10}C_1)^2+2(^{10}C_2)^2+3(^{10}C_3)^2+.....+(^{10}C_{10})^2$ where $^{10}C_r$,r$\in$ {1,2,3,.....,10} denote binomial coefficient then ,the value of $\frac{1}{1430}$X is.......
\item suppose det$\myvec{\sum_{k=0}^{n}k&\sum_{k=0}^{n} {n}C_k k^2\\\sum_{k=0}^{n} {n}C_k k&\sum_{k=0}^{n} {n}C_k 3^k}=0$ holds for some positive integer n.The $\sum_{k=0}^{n} \frac{^{n}C_k}{k+1}$ equals
\item the coefficient of $x^p$ and $x^q$ in the expansion of $(1+x)^{p+q}$ are
\begin{enumerate}
\item equal
\item equal with opposite signs
\item reciprocals of each other
\item none of these
\end {enumerate}
\item If sum of the coefficients in the expansion of $(a+b)^n$ is 4096, then the greatest coefficient in the expansion is 
\begin{enumerate}
\item 1594
\item 792
\item 924
\item 2924
\end{enumerate}
\item the positive integer just greater than \\$(1+0.0001)^{10000} $is 
\begin{enumerate}
\item 4
\item 5
\item 2
\item 3
\end{enumerate}
\item r and n are the positive integers r$>$1,n$>$2 and coefficient of $(r+2)^{th}$ term and $3r^{th}$ tern in the expansion of $(1+x)^{2n}$ are equal,then n equals
\begin{enumerate}
\item 3r
\item 3r+1
\item 2r
\item 2r+1
\end{enumerate}
\item If $\sqrt{7+\sqrt{7+\sqrt{7+.....}}}$ having n radical sign then by methods of mathematical induction which is true
\begin{enumerate}
    \item $a_n>7 \forall n\geq1$
    \item $a_n<7 \forall n\geq1$
    \item $a_n<4 \forall n\geq1$
    \item $a_n<3 \forall n\geq1$
\end{enumerate}
\item If x is positive,the first negative term in the expansion of $(1+x)^\frac{27}{5}$ is 
\begin{enumerate}
\item 6th term
\item 7th term 
\item 5th term 
\item 8th term
\end{enumerate}
\item The number of integral terms in the expansion of $(\sqrt{3}+\sqrt[8]{5})^{256}$ is 
\begin{enumerate}
    \item 35
    \item 32
    \item 33
    \item 34
\end{enumerate}
\item Let S(K)=1+3+5+....+(2K-1)=3+$K^2$ Then which of the following is true
\begin{enumerate}
    \item principle of mathematical induction can be used to prove the formula
    \item S(K)$\rightarrow$S(K+1)
    \item S(K)$\rightarrow$S(K+1)
    \item S(1) is correct
\end{enumerate}
\item The coefficient of middle term in the binomial expansion in powers of x of$(1+\alpha x)^4$ and of $(1-\alpha x)^6$ is the same if $\alpha$ equals to
\begin{enumerate}
    \item $\frac{3}{5}$
    \item $\frac{10}{3}$
    \item $\frac{-3}{10}$
    \item $\frac{-5}{3}$
\end{enumerate}
\item The coefficient of $x^n$ in the expansion of (1+x) $(1-x)^n$ is
\begin{enumerate}
    \item $(-1)^{n-1} n$
    \item $(-1)^n(1-n)$
    \item $(-1)^{n-1}(n-1)^2$
    \item n-1
\end{enumerate}
\item The value of $^{50}C_4+\sum_{r=1}^6$ $^{56-r}C_3$ is
\begin{enumerate}
    \item $^{55}C_4$
    \item $^{55}C_3$
    \item $^{56}C_3$
    \item $^{56}C_4$
\end{enumerate}
\item If A=\myvec{1&0\\1&1} and I=\myvec{1&0\\0&1}, then which one of the following holds for all n $\geq$  1, by the principle of mathematical induction
\begin{enumerate}
    \item $A^n=nA-(n-1)I$
    \item $A^n=2^{n-1}A-(n-1)I$
    \item $A^n=nA+(n-1)I$
    \item $A^n=2^{n-1}A+(n-1)I$
\end{enumerate}
\item If the coefficient of $x^7$ in $[ax^2+\frac{1}{bx}]^{11}$ equals the coefficient of $x^{-7}$ in $[ax-\frac{1}{bx^2}]^{11}$, then a and b satisfy the relation
\begin{enumerate}
    \item a-b=1
    \item a+b=1
    \item $\frac{a}{b}=1$
    \item ab=1
\end{enumerate}
\item If x is so small that $x^3$ and higher powers of x may be neglected, then $\frac{(1+x)^\frac{3}{2}-(1+\frac{1}{2}x)^3}{(1-x)^\frac{1}{2}}$ may be approximated as 
\begin{enumerate}
    \item $1-\frac{3}{8}x^2$
    \item $3x+\frac{3}{8}x^2$
    \item $-\frac{3}{8}x^2$
    \item $\frac{x}{2}-\frac{3}{8}x^2$
\end{enumerate}
\item If expansion in power of x of the function $\frac{1}{(1-ax)(1-bx)}$ is $a_0+a_1x+a_2 x^2+.....$ then $a_n$ is
\begin{enumerate}
    \item $\frac{b^n-a^n}{b-a}$
    \item $\frac{a^n-b^n}{b-a}$
    \item $\frac{a^{n+1}-b^{n+1}}{b-a}$
    \item $\frac{b^{n+1}-a^{n+1}}{b-a}$
\end{enumerate}
\item For natural numbers m,n if 
\begin{align} 
(1-y)^m (1+y)^n=1+a_1y+a_2y^2+.......
\end{align} and $a_1=a_2=10$, then $\myvec{m\\n}$ is
\begin{enumerate}
\item \myvec{20\\45}
\item \myvec{35\\20}
\item \myvec{45\\35}
\item \myvec{35\\45}
\end{enumerate}
\item In the binomial expansion of $(a-b)^n$, $n \geq 5$, the sum of 5th and 6th terms is zero then a/b equals
\begin{enumerate}
    \item $\frac{n-5}{6}$
    \item $\frac{n-4}{5}$
    \item $\frac{5}{n-4}$
    \item $\frac{6}{n-5}$
\end{enumerate}
\item The sum of the series\\
$^{20}C_0-^{20}C_1+^{20}C_2-^{20}C_3+...-.....+^{20}C_{10}$ is
\begin{enumerate}
    \item 0
    \item $^{20}C_{10}$
    \item $-^{20}C_{10}$
    \item $\frac{1}{2}^{20}C_{10}$
\end{enumerate}
\item \textbf{statement-1}:\\$\sum_{r=0}^{n}(r+1) ^nC_r=(n+2)2^{n-1}$\\
\textbf{statement-2}:\\$\sum_{r=0}^{n}(r+1) ^nC_r x^r=(1+x)^n+nx(1+x)^{n-1}$
\begin{enumerate}
    \item statement-1 is true,statement-2 is true
    \item statement-1 is true,statement-2 is true;statement-2 is correct explanation for statement-1
    \item statement-1 is true,statement-2 is true;statement-2 is not correct explanation for statement-1
    \item statement-1 is true,statement-2 is false
\end{enumerate}
\item The reminder left out when $8^{2n}-(62)^{2n+1}$ is divided by 9 is:
\begin{enumerate}
    \item 2
    \item 7
    \item 8
    \item 0
\end{enumerate}
\item $S_1=\sum_{j=1}^{10} j(j-1) ^{10}C_J,\\
S_2=\sum_{j=1} ^{10}$ $^{10}C_j$ and \\
$S_3=\sum_{j=1}^{10} j^2$ $^{10}C_j$\\
\textbf{statement-1}:$S_3=55\times2^9$\\
\textbf{statement-2}:$S_1=90\times2^8$ and $S_2=10\times2^8$
\begin{enumerate}
    \item statement-1 is true,statement-2 is true;statement-2 is not correct explanation for statement-1
    \item statement-1 is true,statement-2 is false
    \item statement-1 is false,statement-2 is true
    \item statement-1 is true,statement-2 is true;statement-2 is  correct explanation for statement-1
\end{enumerate}
\item The Coefficient of $x^7$ in the expansion of\\
 $(1-x-x^2+x^3)^6$ is
\begin{enumerate}
    \item -132
    \item -144
    \item 132
    \item 144
\end{enumerate}
\item If n is a positive integer, then\\
$(\sqrt{3}+1)^{2n}-(\sqrt{3}-1)^{2n}$ is:
\begin{enumerate}
    \item an irrational number 
    \item an odd positive integer
    \item an even positive integer
    \item a rational number other than positive integer
\end{enumerate}
\item The term independent of x in the expansion $(\frac{x+1}{x^\frac{2}{3}-x^\frac{1}{3}+1}-\frac{x-1}{x-x^\frac{1}{2}})^{10}$ is
\begin{enumerate}
    \item 4
    \item 120
    \item 210
    \item 310
\end{enumerate}
\item If the coefficient of $x^3$ and $x^4$ in the expansion of $(1+ax+bx^2)(1-2x)^{18}$ in powers of x both zero,then $\myvec{a\\b}$ is equal to:
\begin{enumerate}
    \item $\myvec{14\\\frac{272}{3}}$
    \item $\myvec{16\\\frac{272}{3}}$
    \item $\myvec{16\\\frac{251}{3}}$
    \item $\myvec{14\\\frac{251}{3}}$
\end{enumerate}
\item The sum of the coefficients of integral power of x in the binomial expansion \\
$(1-2\sqrt{x})^{50}$ is:
\begin{enumerate}
    \item $\frac{1}{2} (3^{50}-1)$
    \item $\frac{1}{2} (2^{50}+1)$
    \item $\frac{1}{2} (3^{50}+1)$
    \item $\frac{1}{2} (3^{50})$
\end{enumerate}
\item The number of terms in the expansion of \\$(1-\frac{2}{x}+\frac{4}{x^2})^n$, $x \neq 0$, is 28, Then the sum of the all the terms in the expansion is:
\begin{enumerate}
    \item 243
    \item 729
    \item 64
    \item 2187
\end{enumerate}
\item The value of $(^{21}C_1- ^{10}C_1)+(^{21}C_2- ^{10}C_2)+(^{21}C_3- ^{10}C_3)+(^{21}C_4- ^{10}C_4)+...+(^{21}C_{10}- ^{10}C_{10})$ is
\begin{enumerate}
    \item $2^{20}-2^{10}$
    \item $2^{21}-2^{11}$
    \item $2^{21}-2^{10}$
    \item $2^{20}-2^{9}$
\end{enumerate}
\item The sum of the all coefficients of all odd degree terms in the expansion of \\$(x+\sqrt{x^3-1})^5+(x-\sqrt{x^3-1})^5$, (x$>$1) is:
\begin{enumerate}
    \item 0
    \item 1
    \item 2
    \item -1
\end{enumerate}
\item If the forth term in the Binomial expansion of $(\frac{2}{x}+x^{log 8x})^6$ (x$>$0) is 20$\times$8$^7,$ then a value of x is:
\begin{enumerate}
    \item $8^3$
    \item $8^2$
    \item $8$
    \item $8^{-2}$
\end{enumerate}
\end{enumerate}
%\end{document}
    

\section{JEE Exercises:Permutation and Combination}
\renewcommand{\theequation}{\theenumi}
\begin{enumerate}[label=\arabic*.,ref=\thesubsection.\theenumi]
\numberwithin{equation}{enumi}

\item In a certain test, $a_i$ students gave wrong answers to atleast i questions, where $i = 1,2,...,k.$ No student gave more than k wrong answers. The total number of wrong answers given is.....\\
\item The side AB, BC and CA of a triangle ABC have 3, 4 and 5 interior points respectively on them. The number of triangles that can be constructed using these interior points as vertices is.....\\
\item Total number of ways in which six $'+'$ and four $'-'$ signs can be arranged in a line such that no two $'-'$ signs occur together is......\\
\item There are four balls of different colours and four boxes of colours, same as those of the balls. The number of ways in which the balls, one each in a box, could be placed such that a ball does not go to a box of its own colour is.....\\
\item The product of any r consecutive natural numbers is always divisible by r!\\
\item $\nCr{n}{r-1} = 36$, $\nCr{n}{r} = 84$ and $\nCr{n}{r+1} = 126$, then r is :
\begin{enumerate}
\item 1
\item 2
\item 3
\item None of these.\\
\end{enumerate}
\item Ten different letters of an alphabet are given. Words with five letters are formed from these given letters. Then the number of words which have at least one letter repeated are
\begin{enumerate}
\item 69760
\item 30240
\item 99748
\item none of these\\
\end{enumerate}
\item The value of the expression $\nCr{47}{4}+\sum_{j=1}^{5} \nCr{52-j}{3}$ is equal to
\begin{enumerate}
\item $\nCr{47}{5}$
\item $\nCr{52}{5}$
\item $\nCr{52}{4}$
\item none of these\\
\end{enumerate} 
\item Eight chairs are numbered 1 to 8. Two women and three men wish to occupy one chair each. First the women choose the chairs from amongst the chairs marked 1 to 4; and then the men select the chairs from amongst the remaining. The number of possible arrangments is
\begin{enumerate}
\item $\nCr{6}{3}$ $\times$ $\nCr{4}{2}$
\item $\nPr{4}{2}$ $\times$ $\nPr{4}{3}$
\item $\nCr{4}{2}$ + $\nPr{4}{3}$
\item none of these\\
\end{enumerate}
\item A five-digit numbers divisible by 3 is to be formed using the numerals 0, 1, 2, 3, 4 and 5, without repetition. The total number of ways this can be done is
\begin{enumerate}
\item 216
\item 240
\item 600
\item 3125\\
\end{enumerate}
\item How many different nine digit numbers can be formed from the number 223355888 by rearranging its digits so that the odd digits occupy even positions?
\begin{enumerate}
\item 16
\item 36
\item 60
\item 180\\
\end{enumerate}
\item Let $T_n$ denote the number of triangles which can be formed using the vertices of a regular polygon of n sides. If $T_{n+1}-T_n = 21$, then n equals
\begin{enumerate}
\item 5
\item 7
\item 6
\item 4\\
\end{enumerate}
\item The number of arrangements of the letters of the word BANANA in which the two N's do not apppear adjacently is
\begin{enumerate}
\item 40
\item 60
\item 80
\item 100\\
\end{enumerate}
\item A rectangle with sides of length (2m-1) and (2n-1) units is divided into squares of unit length by drawing parallel lines as shown in the diagram, then the number of rectangles possible with odd side lengths is

%\includegraphics[width=30mm,scale=0.5]{/home/nlap23679/Downloads/squares of unit length.jpg} 
\begin{enumerate}
\item $(m+n-1)^2$
\item $4^{m+n-1}$
\item $m^2n^2$
\item $m(m+1)n(n+1)$\\
\end{enumerate}
\item If the LCM of p, q is $r^2t^4s^2$, where r, s, t are prime numbers and p,q are the positive integers then the number of ordered pair $(p,q)$ is
\begin{enumerate}
\item 252
\item 254
\item 225
\item 224\\
\end{enumerate}
\item The letters of the word COCHIN are permuted and all the permutations are arranged in an alphabetical order as in an English dictionary. The number of words that appear before the word COCHIN is 
\begin{enumerate}
\item 360
\item 192
\item 96
\item 48\\
\end{enumerate}
\item The number of seven digit integers, with sum of the digits equal to 10 and formed by using the digits 1,2 and 3 only, is
\begin{enumerate}
\item 55
\item 66
\item 77
\item 88\\
\end{enumerate}
\item The total number of ways in which 5 balls of different colours can be distributed among 3 persons so that each person gets at least one ball is 
\begin{enumerate}
\item 75
\item 150
\item 210
\item 243\\
\end{enumerate}
\item Six cards and six envelopes are numbered 1,2,3,4,5,6 and cards are to be placed in envelopes so that each envelope contains exactly one card and no card is placed in the envelope bearing the same number and moreover the card numbered 1 is always placed in envelope numbered 2. Then the number of ways it can be done is 
\begin{enumerate}
\item 264
\item 265
\item 53
\item 67\\
\end{enumerate} 
\item A Debate club consists of 6 girls and 4 boys. A team of 4 members is to be selected from this club including the selection of a captain (from among these 4 members) for the team. If the team has to include  at most one boy, then the number of ways of selecting the team is 
\begin{enumerate}
\item 380
\item 320
\item 260
\item 95\\
\end{enumerate}
\item An n-digit number is a positive number with exactly n digits. Nine hundred distinct n-digit numbers are to be formed using only three digits 2,5 and 7. The smallest value of n for which this is possible is
\begin{enumerate}
\item 6
\item 7
\item 8
\item 9\\
\end{enumerate}
\item Six X's have to be placed in the sqaures of figure given below in such a way that each row contains atleast one X. In how many different ways can this be done
%\includegraphics[width=30mm,scale=0.5]{/home/nlap23679/Downloads/six squares.jpg} \\
\item Five balls of different colours are to be placed in the boxes of different size. Each box can hold all five. In how many different ways can we place the balls so that no box remains empty?\\
\item m men and n women are to be seated in a row so that no two women sit together. If $m>n$, then show that the number of ways in which they can be seated is $\frac{m!(m+1)!}{(m-n+1)!}$\\
\item 7 relatives of a man comprises 4 ladies and 3 gentlemen; his wife also has 7 relatives; 3 of them are ladies and 4 gentlemen. In how many ways they can invite a dinner party of 3 ladies and 3 gentlemen so that there are 3 of man's relatives and 3 of wife's relatives?\\
\item A box contains two white balls, three black balls, four red balls. In how many ways can three balls be drawn from the box if at least one black ball is to be included in the draw?\\
\item Eighteen guests have to be seated, half on each side of a long table. Four particular guests desire to sit on one particular side and three others on the others side. Determine the number of ways in which the sitting arragements can be made.\\
\item A committee of 12 is to be formed from 9 women and 8 men.In how many ways this can be done if at least five women have to be included in a committee? In how many of these committees
\begin{enumerate}
\item The women are in majority?
\item The men are in majority?\\
\end{enumerate}
\item Prove by permutation or otherwise $\frac{(n^2)!}{(n!)^n}$ is an integer $(n\epsilon I^+)$.\\
\item If total number of runs scored in n matches is $(\frac{n+1}{4})$ $(2^{n+1}-n-2)$ where $n>1$, and the runs scored in the kth match are given by K. $2^{n+1-k}$, where $1\leq k\leq n.$ Find n.\\

Let $a_n$ denote the number of all n-digit positive integers formed by the digits 0, 1 or both such that no consecutive digits in them are 0. Let $b_n=$ the number of such n-digit integers ending with digit 1 and $c_n=$ the number of such n-digit integers ending with digit 0.\\
\item The value of $b_6$ is
\begin{enumerate}
\item 7
\item 8
\item 9
\item 11\\
\end{enumerate}
\item which of the following is correct?
\begin{enumerate}
\item $a_{17} = a_{16} + a_{15}$
\item $c_{17} \neq c_{16} + c_{15}$
\item $b_{17} \neq b_{16} + c_{16}$
\item $a_{17} = c_{17} + b_{16}$\\
\end{enumerate}
\item Consider the set of eight vectors V=$\lbrace$a\^{i}+b\^{j}+c\^{k}:a,b,c$\epsilon\lbrace-1,1\rbrace$ $\rbrace.$ Three non-coplanar vectors can be chosen from V in $2^p$ ways. Then p is\\

\item Let $n_1<n_2<n_3<n_4<n_5$ be positive integers such that $n_1+n_2+n_3+n_4+n_5=20.$ Then the number of such distinct arrangements $(n_1,n_2,n_3,n_4,n_5)$ is\\
\item Let $n\geq2$ be an integer. Take n distinct points on a circle and join each pair of points by a line segment. Colour the line segment by joining every pair of adjacent points by blue and the rest by red. If the number of red and blue line segments are equal, then the value of n is\\
\item Let n be number of ways in which 5 boys and 5 girls can stand in a queue in such a way that all the girls stand consecutively in the queue. Let m be the number of ways in which 5 boys and 5 girls can stand in a queue in such a way that exactly four girls stand consecutively in the queue. Then the value of $\frac{m}{n}$ is\\
\item Words of length 10 are formed using the letters A, B, C, D, E, F, G, H, I, J. Let x be the number of such words where no letter is repeated; and let y be the number of such words where exactly one letter is repeated twice and no other letter is repeated. Then, $\frac{y}{9x} = $\\
\item The number of 5 digit numbers which are divisible by 4, with digits from the set $\left\lbrace1,2,3,4,5\right\rbrace$ and the repetition of digits is allowed, is\\
\item Let $\vert X \vert$ denote the number of elements in a set X. Let S$=\left\lbrace1,2,3,4,5,6\right\rbrace$ be a sample space, where each element is likely to occur. If A and B are independent events associated with S, then the number of ordered pairs (A,B) such that $1\leq\vert B \vert < \vert A \vert$, equals.\\
\item Five persons A, B, C, D and E are seated in a circular arrangement. If each of them is given a hat of one of the three colours red, blue and green, then the number of ways of distributing the hats such that the persons seated in adjacent seats get different coloured hats is...\\ 
\item Total no of four digit odd numbers that can be formed using 0,1,2,3,5,7 (using the repetition allowed) are
\begin{enumerate}
\item 216
\item 375
\item 400
\item 720
\end{enumerate}
\item Number greater than 1000 but less than 4000 is formed using the digits 0,1,2,3,4 (repetition allowed). Their number is 
\begin{enumerate}
\item 125
\item 105 
\item 375
\item 625\\
\end{enumerate}
\item Five digit number divisible by 3 is formed using 0,1,2,3,4 and 5 without repetition. Total number of such numbers are
\begin{enumerate}
\item 312
\item 3125 
\item 120
\item 216\\
\end{enumerate}
\item The sum of integers from 1 to 100 that are divisible by 2 or 5 is
\begin{enumerate}
\item 3000
\item 3050
\item 3600
\item 3250\\
\end{enumerate}
\item If $\nCr{n}{r}$ denotes the number of combination of n things taken r at a time, then the expression $\nCr{n}{r+1}+\nCr{n}{r-1}+2 \times \nCr{n}{r}$ equals
\begin{enumerate}
\item $\nCr{n+1}{r+1}$
\item $\nCr{n+2}{r}$
\item $\nCr{n+2}{r+1}$
\item $\nCr{n+1}{r}$\\
\end{enumerate}
\item A student is to answer 10 out of 13 questions in an examination such that he must choose at least 4 from the first five questions. The number of choices available to him is
\begin{enumerate}
\item 346
\item 140 
\item 196
\item 280\\
\end{enumerate}
\item The number of ways in which 6 men and 5 women can dine at a round table if no two women are to sit together is given by
\begin{enumerate}
\item $7! \times 5!$
\item $6! \times 5!$
\item 30!
\item $5! \times 4!$\\
\end{enumerate}
\item How many ways are there to arrange the letters in the word GARDEN with vowels in alphabetical order
\begin{enumerate}
\item 480
\item 240 
\item 360
\item 120\\
\end{enumerate}
\item The number of ways of distributing 8 identical balls in 3 distinct boxes so that none of the boxes is empty is 
\begin{enumerate}
\item $\nCr{8}{3}$
\item 21
\item $3^8$
\item 5\\
\end{enumerate}
\item If the letters of the word SACHIN are arranged in all possible ways and these words are written out as in dictionary, then the word SACHIN appears at serial number
\begin{enumerate}
\item 601
\item 600 
\item 603
\item 602\\
\end{enumerate}
\item At an election, a voter may vote for any number of candidates, not greater than the number to be elected. There are 10 candidates and 4 are to be selected, if a voter votes for at least one candidate, then the number of ways in which he can vote is
\begin{enumerate}
\item 5040
\item 6210 
\item 385
\item 1110\\
\end{enumerate} 
\item The set S $= \left\lbrace1,2,3.....12\right\rbrace$ is to be partitioned into three sets A, B, C of equal size. Thus $A \cup B \cup C = S$, $A \cap B=B \cap C=A \cap C= \phi$. The number of ways to partition S is 
\begin{enumerate}
\item $\frac{12!}{(4!)^3}$
\item $\frac{12!}{(4!)^4}$
\item $\frac{12!}{3!(4!)^3}$
\item $\frac{12!}{3!(4!)^4}$\\
\end{enumerate} 
\item In a Shop there are five types of icecreams available. A child buys six ice-creams.
Statement-1 : The number of different ways the child can buy the six ice-creams is $\nCr{10}{5}$.
Statement-2 : The number of different ways the child can but the six ice-creams is equal to the number of different ways of arranging 6 A's and 4 B's in a row.
\begin{enumerate}
\item Statement-1 is False, Statement-2 is True
\item Statement-1 is True, Statement-2 is True; Statement-2 is a correct explanation for Statement-1
\item Statement-1 is True, Statement-2 is True; Statement-2 is not a correct explanation for Statement-1
\item Statement-1 is True, Statement-2 is False\\ 
\end{enumerate}
\item How many different words can be formed by jumbling the letters in the word MISSISSIPPI in which no two S are adjacent?
\begin{enumerate}
\item $8.\nCr{6}{4}.\nCr{7}{4}$
\item $6.7.\nCr{8}{4}$
\item $6.8.\nCr{7}{4}$
\item $7.\nCr{6}{4}.\nCr{8}{4}$\\
\end{enumerate}   
\item From 6 different novels and 3 different ditionaries, 4 novels and 1 dictionary are to be selected and arranged in a row on a shelf so that the dictionary is always in the middle. Then the number of such arrangement is :
\begin{enumerate}
\item at least 500 but less than 750
\item at least 750 but less than 1000
\item at least 1000
\item at least 500\\ 
\end{enumerate}
\item There are two urns. Urn A has 3 distinct red balls and Urn B has 9 distinct blue balls. From each urn two balls are taken out at random and then transferred to the other. The number of ways in which this can be done is
\begin{enumerate}
\item 36
\item 66 
\item 108
\item 3\\
\end{enumerate}
\item Statement-1 :  The number of ways of distributing 10 identical balls in 4 distinct boxes such that no box is empty is $\nCr{9}{3}$.
Statement-2 : The number of ways of choosing any 3 places from 9 different places is $\nCr{9}{3}$.
\begin{enumerate}
\item Statement-1 is True, Statement-2 is True; Statement-2 is not a correct explanation for Statement-1
\item Statement-1 is True, Statement-2 is False
\item Statement-1 is False, Statement-2 is True
\item Statement-1 is True, Statement-2 is True; Statement-2 is a correct explanation for Statement-1\\
\end{enumerate} 
\item These are 10 points in a plane, out of these 6 are collinear, if N is the number of triangles formed by joining these points. then:
\begin{enumerate}
\item $n \leq 100$
\item $100 < n \leq 140$
\item $140 < n \leq 190$
\item $n > 190$\\
\end{enumerate}  
\item Assuming the balls to be identical except for difference in colours, the number of ways in which one  or more balls can be selected from 10 white, 9 green, and 7 black balls is :
\begin{enumerate}
\item 880
\item 629
\item 630
\item 879\\
\end{enumerate}
\item Let $T_n$ be the number of all possible traingles formed by joining vertices of an n-sided regular polygon. If $T_{n+1}-T_n = 10$, then the value of n is:
\begin{enumerate}
\item 7
\item 5
\item 10
\item 8\\
\end{enumerate}
\item The number of integers greater than 6000 that can be formed, using the digits 3, 5, 6, 7 and 8, without repetition, is :
\begin{enumerate}
\item 120
\item 72
\item 216
\item 192\\
\end{enumerate}
\item If all the words (with or without meaning) having five letters, formed using the letters of the word SMALL and arranged as in a dictionary; then the position of the word SMALL is :
\begin{enumerate}
\item $52^{nd}$
\item $58^{th}$
\item $46^{th}$
\item $59^{th}$\\
\end{enumerate}
\item A man X has 7 friends, 4 of them are ladies and 3 are men. His wife Y also has 7 friends, 3 of them are ladies and 4 are men. Assume X and Y have no common friends. Then the total number of ways in which X and Y together can throw a party inviting 3 ladies and 3 men, so that 3 friends of each of X and Y are in this party, is :
\begin{enumerate}
\item 484
\item 485
\item 468
\item 469\\
\end{enumerate}
\item From 6 different novels and 3 different dictionaries; 4 novels and 1 dictionary are to be selected and arranged in a row on a shelf so that the dictionary is always in the middle. The number of such arrangements is :
\begin{enumerate}
\item less than 500
\item at least 500 but less than 750
\item at least 750 but less than 1000
\item at least 1000\\ 
\end{enumerate}
\item Consider a class of 5 girls and 7 boys. The number of different teams consisting of 2 girls and 3 boys that can be formed from this class, if there are two specific boys A and B, who refuse to be the members of the same team, is:
\begin{enumerate}
\item 500
\item 200
\item 300
\item 350\\
\end{enumerate} 
\item A committee of 11 members is to be formed from 8 males and 5 females. If m is the number of ways the committee is formed with at least 6 males and n is the number of ways the committee is formed with at least 3 females, then:
\begin{enumerate}
\item m + n = 68
\item m = n = 78
\item n = m - 8
\item m = n = 68\\
\end{enumerate}


Each question contains statements given in two columns, which have to be matched. The statements in Column-I are labelled A, B, C, D, while the statements in Column-II are labelled p, q, r, s and t. Any given statement in Column-I can have correct matching with ONE OR MORE statements in Column-II. The appropriate bubbles corresponding to the answers to these questions have to be darkened as illustrated in the following example :
If the correct matches are A-p, s and t; B-q and r; C-p and q; and D-s, then the correct darkening of bubbles will look like

%\includegraphics[width=30mm,scale=0.5]{/home/nlap23679/Downloads/omr.jpg}\\ 

\item Consider all possible permutations of the letters of the word ENDEANOEL. Match the Statements/Expressions in Column-I with the Statements/Expressions in column-II and indicate your answer by darkening the appropriate bubbles in the $4 \times 4$ matrix given in the ORS.\\
\resizebox{\columnwidth}{14pt}{%
\begin{tabular}{cc}
Column-I &Column-II\\
(A) The number of permutations containing the word ENDEA is&(p)5!\\
(B) The number of permutations in which the letter E occurs in the first and the last positions is&(q)2 $\times$ 5!\\ 
(C) The number of permutations in which none of the letters D, L, N occurs in the last five positions is&(r)7 $\times$ 5!\\
(D) The number of permutations in which the letters A, E, O occur only in odd positions is&(s)21 $\times$ 5!\\
\end{tabular}}\\


\item In a high school, a committe has to be formed from a group of 6 boys $M_1,M_2,M_3,M_4,M_5,M_6$ and 5 girls $G_1,G_2,G_3,G_4,G_5.$
\begin{enumerate}
\item Let $\alpha_1$ be the total number of ways in which the committe can be formed such that the committe has 5 members, having exactly 3 boys and 2 girls.
\item Let $\alpha_2$ be the total number of ways in which the committe can be formed such that the committe has at least 2 members and having an equal number of boys and girls.
\item Let $\alpha_3$ be the total number of ways in which the committe can be formed such that the committe has 5 members at least 2 of them being girls.
\item Let $\alpha_4$ be the total number of ways in which the committe can be formed such that the committe has 4 members, having at least 2 girls and such that both $M_1$ and $G_1$ are NOT in the committe together.
\end{enumerate}
\begin{table}[h!]
\centering
\begin{tabular}{c c} 
 LIST-I & LIST-II\\ [0.5ex] 
 P. The value of $\alpha_1$ is & 1. 136 \\ 
 Q. The value of $\alpha_2$ is & 2. 189\\
 R. The value of $\alpha_3$ is & 3. 192\\
 S. The value of $\alpha_4$ is & 4. 200\\
 & 5. 381\\
 & 6. 461\\ [1ex] 
\end{tabular}
\end{table}
The correct option is:\\
\begin{enumerate}
\item P$\rightarrow$4; Q$\rightarrow$6; R$\rightarrow$2; S$\rightarrow$1
\item P$\rightarrow$1; Q$\rightarrow$4; R$\rightarrow$2; S$\rightarrow$3
\item P$\rightarrow$4; Q$\rightarrow$6; R$\rightarrow$5; S$\rightarrow$2
\item P$\rightarrow$4; Q$\rightarrow$2; R$\rightarrow$3; S$\rightarrow$1
\end{enumerate}

\end{enumerate}

\end{document}


